\documentclass[aspectratio=169,10pt]{beamer}

% ==========================================
% PREAMBLE AND SETUP
% ==========================================

\usetheme{Madrid}
\usecolortheme{seahorse}
\setbeamertemplate{navigation symbols}{}
\setbeamertemplate{footline}[frame number]

% Essential packages
\usepackage[utf8]{inputenc}
\usepackage[T1]{fontenc}
\usepackage{amsmath,amsthm,amssymb}
\usepackage{mathtools}
\usepackage{physics}
\usepackage{braket}
\usepackage{tikz}
\usetikzlibrary{shapes,arrows,positioning,calc,patterns,decorations.pathreplacing,intersections,fit,backgrounds,quantikz2}
\usepackage{pgfplots}
\pgfplotsset{compat=1.18}
\usepackage{algorithm,algorithmic}
\usepackage{listings}
\usepackage{xcolor}
\usepackage{animate}
\usepackage{hyperref}
\usepackage{fontawesome5}
\usepackage{multicol}
\usepackage{array}
\usepackage{booktabs}
\usepackage{tabularx}

% Custom colors for consistency
\definecolor{qcblue}{RGB}{0,114,178}
\definecolor{qcred}{RGB}{213,94,0}
\definecolor{qcgreen}{RGB}{0,158,115}
\definecolor{qcpurple}{RGB}{148,0,211}
\definecolor{qcyellow}{RGB}{240,228,66}
\definecolor{qcgray}{RGB}{86,86,86}

% Code listing style
\lstset{
   basicstyle=\ttfamily\footnotesize,
   keywordstyle=\color{qcblue}\bfseries,
   commentstyle=\color{qcgreen},
   stringstyle=\color{qcred},
   numbers=left,
   numberstyle=\tiny\color{qcgray},
   breaklines=true,
   frame=single,
   backgroundcolor=\color{gray!5}
}

% Teaching notes (hidden in compiled output)
\newcommand{\teaching}[1]{}

% Slide header
\newcommand{\slideheader}[1]{\frametitle{#1}}

% Two-column content layout
\long\def\contentcols#1#2#3#4{%
  \begin{columns}[T]
    \begin{column}{#1\textwidth}#3\end{column}
    \begin{column}{#2\textwidth}#4\end{column}
  \end{columns}
}

% Colored boxes
\usepackage{tcolorbox}
\newtcolorbox{redBox}{colback=red!5,colframe=red!60!black,fonttitle=\bfseries}
\newtcolorbox{greenBox}{colback=green!5,colframe=green!60!black,fonttitle=\bfseries}
\newtcolorbox{blueBox}{colback=blue!5,colframe=blue!60!black,fonttitle=\bfseries}
\newtcolorbox{yellowBox}{colback=yellow!5,colframe=yellow!60!black,fonttitle=\bfseries}

% Custom commands
\newcommand{\Z}{\mathbb{Z}}
\newcommand{\R}{\mathbb{R}}
\newcommand{\Zq}{\mathbb{Z}_q}
\newcommand{\bigO}{\mathcal{O}}
\newcommand{\poly}{\text{poly}}

% Theorem environments
\setbeamertemplate{theorems}[numbered]
\theoremstyle{plain}
\newtheorem{proposition}[theorem]{Proposition}

% ==========================================
% TITLE MATTER
% ==========================================

\title[Post-Quantum Cryptography]{Foundations of Post-Quantum Cryptography}
\subtitle{A 3-Hour Introduction for Quantum Technologies Students}
\author[J.R. Martínez]{Dr. José Ramón Martínez Saavedra}
\institute[Quside]{Quside Technologies SL}
\date{\today}

\begin{document}

% ==========================================
% TITLE SLIDE
% ==========================================

% Slide 1.1.1: Title Slide
% Content: Title, subtitle, author, institution, visual hook
% Purpose: Set professional tone, establish context
% Build: Use split background with classical lock being shattered by quantum wave
\begin{frame}[plain]
   \begin{tikzpicture}[remember picture,overlay]
       % Background gradient
       \fill[left color=gray!20, right color=qcblue!20] 
           (current page.south west) rectangle (current page.north east);
       
       % Classical lock (left side)
       %\node[opacity=0.3] at ([xshift=-2cm]current page.west) {
       %    \includegraphics[width=5cm]{placeholder-lock} TODO
       %};
       
       % Quantum circuit (right side)
       \node[opacity=0.3] at ([xshift=2cm]current page.east) {
           \begin{quantikz}
               \qw & \gate{H} & \ctrl{1} & \gate{H} & \meter{} \\
               \qw & \qw & \targ{} & \qw & \meter{}
           \end{quantikz}
       };
   \end{tikzpicture}
   
   \vspace{2cm}
   \begin{center}
       {\Huge \textbf{Foundations of Post-Quantum Cryptography}}\\[0.5cm]
       {\Large A 3-Hour Introduction for Quantum Technologies Students}\\[1cm]
       {\large Dr. José Ramón Martínez Saavedra}\\[0.3cm]
       {\large Quside Technologies SL}\\[0.5cm]
       {\large \today}
   \end{center}
   
\end{frame}


% ==========================================
% SECTION 1.1: INTRODUCTION & HOOK
% ==========================================

\section{Introduction \& Hook}

% Slide 1.1.2: The Invisible Infrastructure - Live Demo
% Content: Live browser demo, certificate details, crypto handshake animation
% Purpose: Hook audience with tangible real-time example
% Build: Create live demo slide with space for browser screenshot and annotation
\begin{frame}{The Invisible Infrastructure}
   \begin{columns}
       \begin{column}{0.6\textwidth}
           \begin{block}{Live Demo}
               \begin{enumerate}
                   \item Open \url{https://github.com}
                   \item Click padlock icon
                   \item View certificate details
               \end{enumerate}
           \end{block}
           
           \vspace{0.5cm}
           
           \begin{alertblock}{In 400 milliseconds...}
               \begin{itemize}
                   \item 3 digital signatures verified
                   \item ECDHE key exchange performed
                   \item 6 cryptographic keys derived
                   \item AES-256-GCM encryption started
               \end{itemize}
           \end{alertblock}
       \end{column}
       
       \begin{column}{0.4\textwidth}
           \begin{tikzpicture}[scale=0.8]
               % Simplified TLS handshake
               \node[circle,draw,thick] (client) at (0,4) {You};
               \node[circle,draw,thick] (server) at (0,0) {GitHub};
               
               \draw[->,thick,qcblue] (client) -- node[right] {\small ClientHello} (server);
               \pause
               \draw[<-,thick,qcgreen] (client) -- node[left] {\small ServerHello} (server);
               \pause
               \draw[<-,thick,qcred] (client) -- node[right] {\small Certificate} (server);
               \pause
               \draw[->,thick,qcpurple] (client) -- node[left] {\small Finished} (server);
           \end{tikzpicture}
       \end{column}
   \end{columns}
   
   \vspace{0.5cm}
   \centering
   \textbf{\Large \color{qcred}All of this will be broken by quantum computers!}
\end{frame}

% Slide 1.1.3: The Scale of Cryptographic Dependence
% Content: Infographic showing daily crypto operations
% Purpose: Establish ubiquity and critical importance
% Build: Animated counters showing increasing numbers
\begin{frame}{One Day in the Digital World}
   \begin{center}
       \begin{tikzpicture}[scale=0.9]
           % Central node
           \node[circle, draw, ultra thick, minimum size=2cm, fill=qcblue!20] (center) at (0,0) {
               \textbf{24 Hours}
           };
           
           % Surrounding statistics
           \node[draw, rounded corners, fill=qcgreen!20] (sig) at (-4,3) {
               \textbf{300 Billion}\\Digital Signatures
           };
           \draw[->, thick] (sig) -- (center);
           
           \node[draw, rounded corners, fill=qcred!20] (email) at (4,3) {
               \textbf{200 Million}\\Encrypted Emails
           };
           \draw[->, thick] (email) -- (center);
           
           \node[draw, rounded corners, fill=qcyellow!20] (video) at (-4,-3) {
               \textbf{500 Hours}\\Video Uploaded
           };
           \draw[->, thick] (video) -- (center);
           
           \node[draw, rounded corners, fill=qcpurple!20] (trans) at (4,-3) {
               \textbf{150,000}\\Financial Transactions
           };
           \draw[->, thick] (trans) -- (center);
           
        % \pause
           % Personal scale
           \node[draw, ultra thick, red, rounded corners] at (0,-5) {
               \textbf{Your Phone Alone: 1,000+ Crypto Operations Daily}
           };
       \end{tikzpicture}
   \end{center}
\end{frame}

% Slide 1.1.4: The Quantum Timeline
% Content: Timeline showing data harvesting now, quantum threats later
% Purpose: Create urgency - this isn't a future problem
% Build: Animated timeline with expanding threat zones
\begin{frame}{The Quantum Timeline}
   \begin{tikzpicture}[scale=0.9]
       % Timeline
       \draw[ultra thick, ->] (0,0) -- (12,0) node[right] {Time};
       
       % Year markers
       \foreach \x/\year in {0/2024, 2/2028, 4/2033, 6/2036, 10/2045} {
           \draw[thick] (\x,0.1) -- (\x,-0.1);
           \node[below] at (\x,-0.3) {\year};
       }
       
       % Events
       \pause
       \fill[qcred!50] (0,0.5) rectangle (12,2);
       \node[white] at (6,1.25) {\textbf{Data Harvesting Happening NOW}};
       
       \pause
       \node[draw, thick, fill=white] at (4,3) {IBM: 100k qubits};
       \draw[thick, ->] (4,2.5) -- (4,0.5);
       
       \pause
       \node[draw, thick, fill=white] at (6,3) {Crypto-relevant QC};
       \draw[thick, ->] (6,2.5) -- (6,0.5);
       
       \pause
       % Danger zone
       \draw[ultra thick, qcred, decorate, decoration={snake, amplitude=2mm}] 
           (0,-1.5) -- node[below] {\textbf{Your secrets already at risk!}} (12,-1.5);
   \end{tikzpicture}
\end{frame}

% Slide 1.1.5: Today's Existential Question
% Content: Central question and three sub-questions
% Purpose: Frame the lecture's narrative arc
% Build: Questions appear one by one
\begin{frame}{The Challenge}
   \begin{center}
       {\Huge \color{qcblue}
       How do we protect information\\
       from computers that don't exist yet?
       }
   \end{center}
   
   \vspace{1cm}
   
   \pause
   \begin{columns}
       \begin{column}{0.33\textwidth}
           \centering
           \textbf{\large 1. Foundation}\\[0.5cm]
           \begin{tikzpicture}
               \node[draw, rounded corners, minimum height=2cm, text width=3cm, align=center] {
                   How does current cryptography work?
               };
           \end{tikzpicture}
       \end{column}
       
       \pause
       \begin{column}{0.33\textwidth}
           \centering
           \textbf{\large 2. Threat}\\[0.5cm]
           \begin{tikzpicture}
               \node[draw, rounded corners, minimum height=2cm, text width=3cm, align=center, fill=qcred!20] {
                   Why do quantum computers break it?
               };
           \end{tikzpicture}
       \end{column}
       
       \pause
       \begin{column}{0.33\textwidth}
           \centering
           \textbf{\large 3. Solution}\\[0.5cm]
           \begin{tikzpicture}
               \node[draw, rounded corners, minimum height=2cm, text width=3cm, align=center, fill=qcgreen!20] {
                   What mathematical structures survive?
               };
           \end{tikzpicture}
       \end{column}
   \end{columns}
\end{frame}

% Slide 1.1.6: Why Physicists Will Lead This Revolution
% Content: Venn diagram showing unique physicist position
% Purpose: Empower audience - they have unique qualifications
% Build: Venn diagram components appear sequentially
\begin{frame}{Why YOU Are Uniquely Positioned}
   \begin{center}
       \begin{tikzpicture}[scale=0.8]
           % Three circles
           \def\firstcircle{(0,0) circle (3cm)}
           \def\secondcircle{(3,0) circle (3cm)}
           \def\thirdcircle{(1.5,-2.6) circle (3cm)}
           
           % Draw circles
           \draw[thick, fill=qcblue!20, fill opacity=0.5] \firstcircle;
           \node at (-1.5,1.5) {\textbf{Quantum}};
           \node at (-1.5,1) {\textbf{Mechanics}};
           
           \pause
           \draw[thick, fill=qcred!20, fill opacity=0.5] \secondcircle;
           \node at (4.5,1.5) {\textbf{Classical}};
           \node at (4.5,1) {\textbf{Crypto}};
           
           \pause
           \draw[thick, fill=qcgreen!20, fill opacity=0.5] \thirdcircle;
           \node at (1.5,-4.5) {\textbf{Physics}};
           \node at (1.5,-5) {\textbf{Intuition}};
           
           \pause
           % Intersection
           \node[font=\Large\bfseries] at (1.5,0) {YOU};
           
           % Advantages
           \pause
           \node[draw, rounded corners, align=left] at (7,2) {
               \textbf{Your Advantages:}\\
               • Deep QM understanding\\
               • Lattice experience\\
               • Theory + experiment\\
               • Physical constraints
           };
       \end{tikzpicture}
   \end{center}
\end{frame}

% Slide 1.1.7: Lecture Roadmap
% Content: Three-hour journey visualization
% Purpose: Set expectations and maintain attention
% Build: Path animation showing the journey
\begin{frame}{Our Three-Hour Journey}
   \begin{center}
       \begin{tikzpicture}[scale=0.9]
           % Timeline
           \draw[ultra thick, ->] (0,0) -- (12,0);
           
           % Hour markers
           \foreach \hour/\x in {1/4, 2/8, 3/12} {
               \draw[thick] (\x,0.1) -- (\x,-0.1);
               \node[below] at (\x,-0.3) {Hour \hour};
           }
           
           % Content blocks
           \pause
           \node[draw, thick, rounded corners, fill=qcblue!20, text width=3.5cm, align=center] 
               at (2,2) {\textbf{Foundations}\\Current crypto\\Why it fails};
           \draw[thick, ->] (2,1.2) -- (2,0.1);
           
           \pause
           \node[draw, thick, rounded corners, fill=qcred!20, text width=3.5cm, align=center] 
               at (6,2) {\textbf{Quantum Threat}\\Shor's algorithm\\Lattice solutions};
           \draw[thick, ->] (6,1.2) -- (6,0.1);
           
           \pause
           \node[draw, thick, rounded corners, fill=qcgreen!20, text width=3.5cm, align=center] 
               at (10,2) {\textbf{Solutions}\\NIST standards\\Your future};
           \draw[thick, ->] (10,1.2) -- (10,0.1);
           
           % Interactive elements
           \pause
           \node[star, star points=5, draw, fill=qcyellow!50] at (2,0) {};
           \node[star, star points=5, draw, fill=qcyellow!50] at (6,0) {};
           \node[star, star points=5, draw, fill=qcyellow!50] at (10,0) {};
           \node at (6,-1) {\faIcon{laptop-code} Interactive elements throughout};
       \end{tikzpicture}
   \end{center}
\end{frame}

% ==========================================
% SECTION 1.2: CLASSICAL CRYPTOGRAPHY ESSENTIALS
% ==========================================

\section{Classical Cryptography Essentials}

% Slide 1.2.1: The Three Pillars of Cryptographic Security
% Content: Visual of three pillars holding up digital security
% Purpose: Establish fundamental security requirements
% Build: Pillars appear one by one with icons
\begin{frame}{The Three Pillars of Security}
   \begin{center}
       \begin{tikzpicture}[scale=0.8]
           % Platform
           \draw[ultra thick, fill=gray!30] (-5,-1) rectangle (5,-0.5);
           \node at (0,-0.75) {\textbf{DIGITAL SECURITY}};
           
           % Pillars appear one by one
           \pause
           % Pillar 1
           \draw[thick, fill=qcblue!30] (-3.5,0) rectangle (-2,3);
           \node[rotate=90] at (-2.75,1.5) {\textbf{Confidentiality}};
           \node at (-2.75,3.5) {\faIcon{lock}};
           \node[text width=2cm, align=center] at (-2.75,-2) {Encryption\\(Lock box)};
           
           \pause
           % Pillar 2
           \draw[thick, fill=qcgreen!30] (-0.75,0) rectangle (0.75,3);
           \node[rotate=90] at (0,1.5) {\textbf{Integrity}};
           \node at (0,3.5) {\faIcon{shield-alt}};
           \node[text width=2cm, align=center] at (0,-2) {Hashing\\(Tamper seal)};
           
           \pause
           % Pillar 3
           \draw[thick, fill=qcred!30] (2,0) rectangle (3.5,3);
           \node[rotate=90] at (2.75,1.5) {\textbf{Authentication}};
           \node at (2.75,3.5) {\faIcon{signature}};
           \node[text width=2cm, align=center] at (2.75,-2) {Signatures\\(Identity)};
           
           \pause
           % Warning
           \node[draw, thick, red, rounded corners] at (0,-3.5) {
               \textbf{Remove ANY pillar → Total collapse!}
           };
       \end{tikzpicture}
   \end{center}
\end{frame}

% Slide 1.2.2: Pillar 1 - Confidentiality Through Encryption
% Content: Simple encryption visualization
% Purpose: Define encryption's role and limitations
% Build: Show transformation from plaintext to ciphertext
\begin{frame}{Pillar 1: Confidentiality}
   \begin{columns}
       \begin{column}{0.5\textwidth}
           \begin{center}
               \begin{tikzpicture}[scale=0.8]
                   % Plaintext
                   \node[draw, rounded corners, fill=white] (plain) at (0,3) {
                       \texttt{ATTACK AT DAWN}
                   };
                   
                   \pause
                   % Encryption box
                   \node[draw, thick, fill=qcblue!30, minimum width=3cm, minimum height=1cm] 
                       (enc) at (0,1.5) {\textbf{Encrypt}};
                   \draw[->, thick] (plain) -- (enc);
                   
                   % Key
                   \node[draw, fill=qcyellow!30] (key) at (-3,1.5) {\faIcon{key} Key};
                   \draw[->, thick] (key) -- (enc);
                   
                   \pause
                   % Ciphertext
                   \node[draw, rounded corners, fill=qcgray!30] (cipher) at (0,0) {
                       \texttt{X9kP2m8QzR5n}
                   };
                   \draw[->, thick] (enc) -- (cipher);
               \end{tikzpicture}
           \end{center}
       \end{column}
       
       \begin{column}{0.5\textwidth}
           \begin{block}{What Encryption Provides}
               \begin{itemize}
                   \item Transforms readable → unreadable
                   \item Only key holder can reverse
                   \item Examples: AES, RSA, ChaCha20
               \end{itemize}
           \end{block}
           
           \pause
           \begin{alertblock}{Critical Limitation}
               Encryption alone is NOT enough!
               \begin{itemize}
                   \item No integrity guarantee
                   \item No sender verification
                   \item Vulnerable to tampering
               \end{itemize}
           \end{alertblock}
       \end{column}
   \end{columns}
\end{frame}

% Slide 1.2.3: When Confidentiality Fails
% Content: NSA-Google case study
% Purpose: Show real consequences of encryption failure
% Build: Timeline of the incident
\begin{frame}{Case Study: When Encryption Isn't Used}
   \begin{center}
       \textbf{\Large 2013: NSA vs Google Internal Traffic}
   \end{center}
   
   \begin{tikzpicture}[scale=0.9]
       % Google data centers
       \node[draw, thick, minimum width=2cm, minimum height=1.5cm] (dc1) at (0,0) {
           \includegraphics[width=1cm]{placeholder-server}\\
           \small DC1
       };
       
       \node[draw, thick, minimum width=2cm, minimum height=1.5cm] (dc2) at (6,0) {
           \includegraphics[width=1cm]{placeholder-server}\\
           \small DC2
       };
       
       \pause
       % Internal traffic
       \draw[<->, ultra thick, qcblue] (dc1) -- node[above] {Private fiber} node[below] {\color{red}Unencrypted!} (dc2);
       
       \pause
       % NSA tap
       \node[draw, thick, fill=qcred!30] (nsa) at (3,-2) {NSA Tap};
       \draw[->, thick, red] (3,0) -- (nsa);
       
       \pause
       % Impact
       \node[draw, ultra thick, red, rounded corners] at (3,2) {
           \textbf{Result: Massive privacy breach}
       };
       
       \pause
       % Lesson
       \node[draw, thick, rounded corners, fill=qcgreen!20] at (3,-3.5) {
           \textbf{Lesson: Physical security $\neq$ Cryptographic security}
       };
   \end{tikzpicture}
\end{frame}

% Slide 1.2.4: Pillar 2 - Integrity Through Hashing
% Content: Hash function visualization
% Purpose: Explain integrity verification concept
% Build: Show how small changes create different hashes
\begin{frame}{Pillar 2: Integrity}
   \begin{columns}
       \begin{column}{0.6\textwidth}
           \begin{center}
               \begin{tikzpicture}[scale=0.8]
                   % Original message
                   \node[draw, rounded corners] (msg1) at (0,3) {
                       \texttt{Transfer \$1000 to Alice}
                   };
                   
                   % Hash function
                   \node[draw, thick, fill=qcgreen!30, minimum width=3cm] 
                       (hash1) at (0,1.5) {\textbf{SHA-256}};
                   \draw[->, thick] (msg1) -- (hash1);
                   
                   % Hash output
                   \node[draw, rounded corners, fill=qcgray!20] (out1) at (0,0) {
                       \texttt{7d865e95...}
                   };
                   \draw[->, thick] (hash1) -- (out1);
                   
                   \pause
                   % Modified message
                   \node[draw, rounded corners, fill=qcred!20] (msg2) at (5,3) {
                       \texttt{Transfer \$9000 to Alice}
                   };
                   
                   % Hash function
                   \node[draw, thick, fill=qcgreen!30, minimum width=3cm] 
                       (hash2) at (5,1.5) {\textbf{SHA-256}};
                   \draw[->, thick] (msg2) -- (hash2);
                   
                   % Different hash
                   \node[draw, rounded corners, fill=qcred!30] (out2) at (5,0) {
                       \texttt{2a4b9c12...}
                   };
                   \draw[->, thick] (hash2) -- (out2);
                   
                   \pause
                   % Comparison
                   \draw[<->, ultra thick, red] (out1) -- node[below] {Completely different!} (out2);
               \end{tikzpicture}
           \end{center}
       \end{column}
       
       \begin{column}{0.4\textwidth}
           \begin{block}{Hash Properties}
               \begin{itemize}
                   \item Fixed size output
                   \item Deterministic
                   \item Avalanche effect
                   \item One-way function
               \end{itemize}
           \end{block}
           
           \pause
           \begin{exampleblock}{Use Cases}
               \begin{itemize}
                   \item File verification
                   \item Password storage
                   \item Digital fingerprints
                   \item Blockchain
               \end{itemize}
           \end{exampleblock}
       \end{column}
   \end{columns}
\end{frame}

% Slide 1.2.5: When Integrity Fails
% Content: MD5 collision attack on certificates
% Purpose: Demonstrate critical importance of hash security
% Build: Show how collision creates rogue certificate
\begin{frame}{Case Study: MD5 Collision Disaster (2008)}
   \begin{center}
       \begin{tikzpicture}[scale=0.9]
           % Legitimate certificate
           \node[draw, thick, fill=qcgreen!20, text width=4cm, align=center] (legit) at (-3,2) {
               \textbf{Legitimate Cert}\\
               \small website.com\\
               \small Valid purposes
           };
           
           % Hash
           \node[draw, fill=qcgray!30] (hash1) at (-3,0) {
               MD5: \texttt{3a7f9b2c...}
           };
           \draw[->, thick] (legit) -- (hash1);
           
           \pause
           % Rogue certificate
           \node[draw, thick, fill=qcred!20, text width=4cm, align=center] (rogue) at (3,2) {
               \textbf{Rogue CA Cert}\\
               \small Can sign anything!\\
               \small Ultimate power
           };
           
           % Same hash!
           \node[draw, fill=qcgray!30] (hash2) at (3,0) {
               MD5: \texttt{3a7f9b2c...}
           };
           \draw[->, thick] (rogue) -- (hash2);
           
           \pause
           % Collision
           \draw[<->, ultra thick, red] (hash1) -- node[above] {\textbf{COLLISION!}} (hash2);
           
           \pause
           % Impact
           \node[draw, ultra thick, red, rounded corners, text width=8cm, align=center] at (0,-2) {
               \textbf{Impact: Attackers can impersonate ANY website}\\
               \small Complete breakdown of web security
           };
       \end{tikzpicture}
   \end{center}
\end{frame}

% Slide 1.2.6: Pillar 3 - Authentication Through Signatures
% Content: Digital signature process
% Purpose: Complete the security triad
% Build: Show signing and verification process
\begin{frame}{Pillar 3: Authentication}
   \begin{center}
       \begin{tikzpicture}[scale=0.8]
           % Signing process
           \node[draw, rounded corners] (msg) at (-4,3) {
               \texttt{I authorize \$1M transfer}
           };
           
           \pause
           % Private key
           \node[draw, fill=qcred!30] (privkey) at (-6,1) {
               \faIcon{key} Private Key
           };
           
           % Sign operation
           \node[draw, thick, fill=qcred!30, minimum width=3cm] (sign) at (-4,1) {
               \textbf{Sign}
           };
           \draw[->, thick] (msg) -- (sign);
           \draw[->, thick] (privkey) -- (sign);
           
           \pause
           % Signature
           \node[draw, rounded corners, fill=qcgray!30] (sig) at (-4,-1) {
               \texttt{Digital Signature}
           };
           \draw[->, thick] (sign) -- (sig);
           
           % Verification process
           \pause
           \node[align=center] at (0,1) {\textbf{Anyone can verify}\\with public key};
           
           % Public key
           \node[draw, fill=qcgreen!30] (pubkey) at (2,1) {
               \faIcon{key} Public Key
           };
           
           % Verify operation
           \node[draw, thick, fill=qcgreen!30, minimum width=3cm] (verify) at (4,1) {
               \textbf{Verify}
           };
           \draw[->, thick] (sig) -| (verify);
           \draw[->, thick] (msg) -| (verify);
           \draw[->, thick] (pubkey) -- (verify);
           
           \pause
           % Result
           \node[draw, thick, fill=qcgreen!20] at (4,-1) {
               \faIcon{check-circle} Valid!
           };
           \draw[->, thick] (verify) -- (4,-1);
       \end{tikzpicture}
   \end{center}
   
   \vspace{0.5cm}
   \begin{columns}
       \begin{column}{0.5\textwidth}
           \centering
           \textbf{Properties:}\\
           • Non-repudiation\\
           • Public verifiability\\
           • Unforgeable
       \end{column}
       \begin{column}{0.5\textwidth}
           \centering
           \textbf{Examples:}\\
           • RSA signatures\\
           • ECDSA\\
           • EdDSA
       \end{column}
   \end{columns}
\end{frame}

% Slide 1.2.7: When Authentication Fails
% Content: DigiNotar breach case study
% Purpose: Show authentication failure is catastrophic
% Build: Timeline and impact cascade
\begin{frame}{Case Study: DigiNotar Catastrophe (2011)}
   \begin{columns}
       \begin{column}{0.5\textwidth}
           \begin{center}
               \textbf{\Large Timeline of Disaster}
               
               \begin{tikzpicture}[scale=0.8]
                   % Timeline
                   \draw[thick, ->] (0,0) -- (0,-6);
                   
                   % Events
                   \node[right, align=left] at (0.2,0) {\textbf{June:} Hackers breach DigiNotar};
                   \pause
                   
                   \node[right, align=left] at (0.2,-1.5) {\textbf{July:} Issue *.google.com cert};
                   \pause
                   
                   \node[right, align=left] at (0.2,-3) {\textbf{August:} Iran uses for MITM};
                   \pause
                   
                   \node[right, align=left] at (0.2,-4.5) {\textbf{Sept:} Company bankrupt};
                   \pause
                   
                   \node[right, align=left, red] at (0.2,-6) {\textbf{Result:} Dutch gov crisis};
               \end{tikzpicture}
           \end{center}
       \end{column}
       
       \begin{column}{0.5\textwidth}
           \begin{alertblock}{Impact Cascade}
               \pause
               \begin{enumerate}
                   \item Fraudulent certificates trusted globally
                   \pause
                   \item Iranian citizens' Gmail compromised
                   \pause
                   \item All browsers revoke DigiNotar
                   \pause
                   \item Dutch government sites fail
                   \pause
                   \item Complete infrastructure rebuild
               \end{enumerate}
           \end{alertblock}
           
           \pause
           \begin{block}{Lesson}
               Authentication failure affects everyone who trusts the system
           \end{block}
       \end{column}
   \end{columns}
\end{frame}

% Slide 1.2.8: The Pre-Digital Key Distribution Problem
% Content: Visual showing exponential key growth
% Purpose: Motivate public key cryptography
% Build: Animate growing number of required keys
\begin{frame}{The Original Problem: Key Distribution}
   \begin{center}
       \begin{tikzpicture}[scale=0.9]
           % Small network
           \node[circle, draw, thick] (a) at (0,2) {A};
           \node[circle, draw, thick] (b) at (2,2) {B};
           \node[circle, draw, thick] (c) at (1,0) {C};
           
           % Keys needed
           \draw[thick, qcblue] (a) -- node[above] {\small Key 1} (b);
           \draw[thick, qcred] (b) -- node[right] {\small Key 2} (c);
           \draw[thick, qcgreen] (c) -- node[left] {\small Key 3} (a);
           
           \node at (1,-1.5) {3 people = 3 keys};
           
           \pause
           % Medium network
           \begin{scope}[xshift=5cm]
               \foreach \i in {1,...,5} {
                   \node[circle, draw, thick] (n\i) at ({72*\i}:1.5) {\i};
               }
               \foreach \i in {1,...,5} {
                   \foreach \j in {1,...,5} {
                       \ifnum\i<\j
                           \draw[gray] (n\i) -- (n\j);
                       \fi
                   }
               }
               \node at (0,-2.5) {5 people = 10 keys};
           \end{scope}
           
           \pause
           % Large network indication
           \begin{scope}[xshift=9cm]
               \node[draw, thick, rounded corners, minimum width=3cm, minimum height=3cm, fill=qcred!20] {
                   1000 people\\[0.5em]
                   \textbf{499,500}\\keys!
               };
           \end{scope}
       \end{tikzpicture}
   \end{center}
   
   \vspace{0.5cm}
   \pause
   \begin{alertblock}{The Circular Paradox}
       To communicate securely, you must first communicate securely!
   \end{alertblock}
\end{frame}

% Slide 1.2.9: The 1976 Revolution - Diffie-Hellman Insight
% Content: One-way function with trapdoor visualization
% Purpose: Introduce asymmetric cryptography concept
% Build: Show the revolutionary idea step by step
\begin{frame}{The Revolutionary Insight (1976)}
   \begin{columns}
       \begin{column}{0.5\textwidth}
           \begin{quote}
               ``We stand today on the brink of a revolution in cryptography.''
               
               \vspace{0.5em}
               -- Diffie \& Hellman
           \end{quote}
           
           \vspace{1em}
           \pause
           
           \textbf{The Breakthrough:}\\
           Separate the key into two parts with different capabilities!
           
           \vspace{0.5em}
           \pause
           \begin{itemize}
               \item \textcolor{qcgreen}{Public key: Can encrypt}
               \item \textcolor{qcred}{Private key: Can decrypt}
           \end{itemize}
       \end{column}
       
       \begin{column}{0.5\textwidth}
           \begin{center}
               \begin{tikzpicture}[scale=0.8]
                   % One-way function
                   \node[draw, rounded corners] (x) at (0,3) {Private: $x$};
                   \pause
                   
                   \node[draw, thick, fill=qcblue!30, minimum width=3cm] (f) at (0,1.5) {
                       \textbf{One-way Function}
                   };
                   \draw[->, ultra thick, qcgreen] (x) -- node[right] {Easy} (f);
                   
                   \pause
                   \node[draw, rounded corners] (y) at (0,0) {Public: $f(x)$};
                   \draw[->, ultra thick, qcgreen] (f) -- (y);
                   
                   \pause
                   % Reverse direction
                   \draw[->, ultra thick, red, dashed] (y) to[bend right=60] node[left] {Hard!} (x);
                   
                   \pause
                   % With trapdoor
                   \node[draw, fill=qcyellow!30] (trap) at (3,0) {Trapdoor};
                   \draw[->, thick, qcgreen] (trap) -- (y);
                   \draw[->, thick, qcgreen, dashed] (y) to[bend left=60] node[right] {Easy with trapdoor} (x);
               \end{tikzpicture}
           \end{center}
       \end{column}
   \end{columns}
\end{frame}

% Slide 1.2.10: How Diffie-Hellman Works
% Content: Step-by-step DH key exchange
% Purpose: First concrete example of public key crypto
% Build: Animate the protocol steps
\begin{frame}{Diffie-Hellman Key Exchange}
   \begin{center}
       \begin{tikzpicture}[scale=0.9]
           % Alice and Bob
           \node[circle, draw, thick, fill=qcblue!20] (alice) at (-3,0) {Alice};
           \node[circle, draw, thick, fill=qcred!20] (bob) at (3,0) {Bob};
           
           % Public parameters
           \node[draw, rounded corners, fill=qcgray!20] at (0,3) {
               Public: $g = 5$, $p = 23$
           };
           
           \pause
           % Step 1: Choose secrets
           \node[draw, fill=qcblue!10] at (-3,-1.5) {Secret: $a = 6$};
           \node[draw, fill=qcred!10] at (3,-1.5) {Secret: $b = 15$};
           
           \pause
           % Step 2: Compute public values
           \node[draw, rounded corners] at (-3,-2.5) {$A = 5^6 \bmod 23 = 8$};
           \node[draw, rounded corners] at (3,-2.5) {$B = 5^{15} \bmod 23 = 19$};
           
           \pause
           % Step 3: Exchange
           \draw[->, thick, qcblue] (-2.5,-2.5) to[bend left] node[above] {$A = 8$} (2.5,-2.5);
           \draw[->, thick, qcred] (2.5,-3) to[bend left] node[below] {$B = 19$} (-2.5,-3);
           
           \pause
           % Step 4: Compute shared secret
           \node[draw, rounded corners, fill=qcgreen!20] at (-3,-4) {$K = 19^6 \bmod 23 = 2$};
           \node[draw, rounded corners, fill=qcgreen!20] at (3,-4) {$K = 8^{15} \bmod 23 = 2$};
           
           \pause
           % Eavesdropper
           \node[circle, draw, thick, fill=qcgray!30] (eve) at (0,-1) {Eve};
           \node[align=center] at (0,-2) {Sees: $g=5$, $p=23$\\$A=8$, $B=19$\\
               \color{red}Cannot compute $K$!};
       \end{tikzpicture}
   \end{center}
\end{frame}


% ==================================================
% HOUR 2: QUANTUM THREAT AND LATTICE FOUNDATIONS
% ==================================================

\section{Quantum Algorithms Breaking Crypto}

% Slide 2.1.1: Section Title - The Quantum Cryptographic Apocalypse
\begin{frame}
  \slideheader{The Quantum Cryptographic Apocalypse}
  
  \vfill
  \begin{center}
    \begin{tikzpicture}[scale=1.2]
      % Quantum computer icon
      \node[draw, ultra thick, minimum width=3cm, minimum height=2cm, fill=purple!20] (qc) at (0,2) {
        \Large \textbf{Quantum\\Computer}
      };
      
      % Lightning bolts
      \draw[ultra thick, red!80, decorate, decoration={zigzag, segment length=3mm, amplitude=2mm}] 
        (-0.5,0.5) -- (-2,-1.5);
      \draw[ultra thick, red!80, decorate, decoration={zigzag, segment length=3mm, amplitude=2mm}] 
        (0,0.5) -- (0,-1.5);
      \draw[ultra thick, red!80, decorate, decoration={zigzag, segment length=3mm, amplitude=2mm}] 
        (0.5,0.5) -- (2,-1.5);
      
      % Broken crypto logos
      \node[opacity=0.3, rotate=25] at (-2,-2) {\includegraphics[width=1.5cm]{example-image}};
      \node[opacity=0.3, rotate=-15] at (0,-2) {\includegraphics[width=1.5cm]{example-image}};
      \node[opacity=0.3, rotate=35] at (2,-2) {\includegraphics[width=1.5cm]{example-image}};
      
      % Text overlay
      \node[font=\Large\bfseries, red!80] at (0,-3.5) {RSA};
      \node[font=\Large\bfseries, red!80] at (-2,-3.5) {ECC};
      \node[font=\Large\bfseries, red!80] at (2,-3.5) {DSA};
    \end{tikzpicture}
    
    \vspace{1cm}
    \Huge \textbf{``Not if, but when''}
  \end{center}
  \vfill
  
  \teaching{Reset attention after break with dramatic visuals}
\end{frame}

% Slide 2.1.2: Dispelling the Parallel Search Myth
\begin{frame}
  \slideheader{Dispelling the Parallel Search Myth}
  
  \begin{columns}[T]
    \begin{column}{0.5\textwidth}

    \begin{blueBox}
      \textbf{Common Misconception:}
      \begin{itemize}
        \item ``QC tries all keys at once''
        \item ``Quantum = massive parallelism''
        \item ``Like having $2^n$ classical computers''
      \end{itemize}
    \end{blueBox}
    
    \vspace{0.5cm}
    
    \begin{redBox}
      \textbf{Reality:}
      \begin{itemize}
        \item Quantum interference patterns
        \item Amplitude cancellation/reinforcement
        \item Extracts global properties
      \end{itemize}
    \end{redBox}
  
    \end{column}
    \begin{column}{0.5\textwidth}

    \begin{tikzpicture}[scale=0.8]
      % Superposition visualization
      \draw[thick] (0,0) -- (5,0) node[right] {$x$};
      \draw[thick] (0,-2) -- (0,2) node[above] {Amplitude};
      
      % Wave functions
      \foreach \x in {0.5,1,1.5,2,2.5,3,3.5,4,4.5} {
        \draw[blue!60, thick] (\x,-1.5) -- (\x,1.5);
      }
      
      % Interference pattern
      \draw[red, ultra thick, smooth] plot coordinates {
        (0,0) (0.5,1) (1,0) (1.5,-1) (2,0) (2.5,1.5) (3,0) (3.5,-1.5) (4,0) (4.5,1) (5,0)
      };
      
      \node[below] at (2.5,-2.5) {\textbf{Interference reveals structure}};
    \end{tikzpicture}
    
    \vspace{0.3cm}
    
    \begin{greenBox}
      \textbf{Better Analogy:}\\
      Like a diffraction grating revealing hidden wavelengths
    \end{greenBox}
  
    \end{column}
  \end{columns}
  
  \teaching{Correct fundamental misunderstanding early}
\end{frame}

% Slide 2.1.3: Shor's Algorithm - The Real Magic
\begin{frame}
  \slideheader{Shor's Algorithm - The Real Magic}
  
  \begin{columns}[T]
    \begin{column}{0.55\textwidth}

    \begin{center}
      \textbf{\large Not brute force - finds hidden periods!}
    \end{center}
    
    \vspace{0.3cm}
    
    \begin{enumerate}
      \item \textbf{Transform problem:}\\
        Factoring $\rightarrow$ Period finding
        
      \item \textbf{Create superposition:}\\
        $\sum_x |x\rangle|f(x)\rangle$ over all values
        
      \item \textbf{Extract period:}\\
        QFT reveals through interference
    \end{enumerate}
    
    \vspace{0.5cm}
    
    \begin{greenBox}
      \textbf{Physics Analogy:}\\
      Diffraction grating reveals light wavelength $\lambda$ from interference pattern
    \end{greenBox}
  
    \end{column}
    \begin{column}{0.45\textwidth}

    \begin{tikzpicture}[scale=0.7]
      % Periodic function
      \draw[thick, ->] (0,0) -- (5,0) node[right] {$x$};
      \draw[thick, ->] (0,0) -- (0,3) node[above] {$f(x)$};
      
      % Draw periodic pattern
      \foreach \i in {0,...,4} {
        \draw[blue, ultra thick] plot[smooth] coordinates {
          ({\i},0) ({\i+0.25},2) ({\i+0.5},1) ({\i+0.75},1.5) ({\i+1},0)
        };
      }
      
      % Period indicator
      \draw[<->, red, thick] (0,-0.5) -- (1,-0.5) node[midway, below] {period $r$};
      
      % QFT arrow
      \node[draw, thick, fill=yellow!30] (qft) at (7,1.5) {QFT};
      \draw[->, ultra thick] (5.5,1.5) -- (6.5,1.5);
      
      % Frequency domain
      \draw[thick, ->] (8,0) -- (11,0) node[right] {freq};
      \draw[thick, ->] (8,0) -- (8,3);
      
      % Peak at frequency
      \draw[red, ultra thick] (9,2.5) -- (9,0);
      \node[red, below] at (9,0) {$1/r$};
    \end{tikzpicture}
  
    \end{column}
  \end{columns}
  
  \teaching{Core conceptual understanding - it's about finding patterns}
\end{frame}

% Slide 2.1.4: From Factoring to Period Finding
\begin{frame}
  \slideheader{From Factoring to Period Finding}
  
  \begin{center}
    \textbf{\large Mathematical Connection}
  \end{center}
  
  \begin{blueBox}
    \textbf{Goal:} Factor $n = pq$ (RSA modulus)
  \end{blueBox}
  
  \vspace{0.3cm}
  
  \begin{enumerate}
    \item Choose random $a$ coprime to $n$
    
    \item Define function: $f(x) = a^x \bmod n$
    
    \item \textbf{Key insight:} This function is periodic!\\
    Find period $r$ such that $a^r \equiv 1 \pmod{n}$
    
    \item If $r$ is even and $a^{r/2} \not\equiv -1 \pmod{n}$:
    \begin{align}
      a^r - 1 &\equiv 0 \pmod{n}\\
      (a^{r/2} - 1)(a^{r/2} + 1) &\equiv 0 \pmod{n}
    \end{align}
    
    \item Then $\gcd(a^{r/2} - 1, n)$ gives a factor!
  \end{enumerate}
  
  \vspace{0.3cm}
  
  \begin{greenBox}
    \textbf{Success probability} $\geq 1/2$ for random $a$
  \end{greenBox}
  
  \teaching{Show the mathematical reduction clearly}
\end{frame}

% Slide 2.1.5: Concrete Example - Factoring 15
\begin{frame}
  \slideheader{Concrete Example - Factoring 15}
  
  \begin{columns}[T]
    \begin{column}{0.5\textwidth}

    \textbf{Let's factor $n = 15$ using $a = 7$:}
    
    \vspace{0.3cm}
    
    \begin{tabular}{ll}
      $7^0 \bmod 15$ & $= 1$ \\
      $7^1 \bmod 15$ & $= 7$ \\
      $7^2 \bmod 15$ & $= 49 \bmod 15 = 4$ \\
      $7^3 \bmod 15$ & $= 28 \bmod 15 = 13$ \\
      $7^4 \bmod 15$ & $= 91 \bmod 15 = \mathbf{1}$ \\
    \end{tabular}
    
    \vspace{0.3cm}
    
    \begin{redBox}
      Period found: $r = 4$
    \end{redBox}
    
    \vspace{0.3cm}
    
    Since $r = 4$ is even:\\
    $a^{r/2} - 1 = 7^2 - 1 = 48$
    
    $\gcd(48, 15) = \gcd(48, 15) = 3$
    
    \begin{greenBox}
      $15 = 3 \times 5$ \checkmark
    \end{greenBox}
  
    \end{column}
    \begin{column}{0.5\textwidth}

    \begin{tikzpicture}[scale=0.8]
      % Plot the periodic function
      \draw[thick, ->] (0,0) -- (5,0) node[right] {$x$};
      \draw[thick, ->] (0,0) -- (0,3.5) node[above] {$a^x \bmod n$};
      
      % Plot points
      \node[circle, fill=blue, inner sep=2pt] (p0) at (0,0.2) {};
      \node[circle, fill=blue, inner sep=2pt] (p1) at (1,1.4) {};
      \node[circle, fill=blue, inner sep=2pt] (p2) at (2,0.8) {};
      \node[circle, fill=blue, inner sep=2pt] (p3) at (3,2.6) {};
      \node[circle, fill=blue, inner sep=2pt] (p4) at (4,0.2) {};
      
      % Labels
      \node[below] at (0,0) {$0$};
      \node[below] at (1,0) {$1$};
      \node[below] at (2,0) {$2$};
      \node[below] at (3,0) {$3$};
      \node[below] at (4,0) {$4$};
      
      \node[left] at (0,0.2) {$1$};
      \node[left] at (0,1.4) {$7$};
      \node[left] at (0,0.8) {$4$};
      \node[left] at (0,2.6) {$13$};
      
      % Connect points
      \draw[blue, thick] (p0) -- (p1) -- (p2) -- (p3) -- (p4);
      
      % Highlight period
      \draw[<->, red, ultra thick] (0,-0.5) -- (4,-0.5) node[midway, below] {period = 4};
    \end{tikzpicture}
  
    \end{column}
  \end{columns}
  
  \teaching{Demystify with concrete, small numbers}
\end{frame}

% Slide 2.1.6: Quantum Circuit for Shor
\begin{frame}
  \slideheader{Quantum Circuit for Shor}
  
  \begin{center}
    \textbf{\large High-level Circuit Structure}
  \end{center}
  
  \begin{tikzpicture}[scale=0.9]
    % Input registers
    \node[left] at (-1,3) {$|0\rangle^{\otimes n}$};
    \node[left] at (-1,1) {$|0\rangle^{\otimes n}$};
    
    % Quantum wires
    \draw[thick] (0,3) -- (12,3);
    \draw[thick] (0,2.8) -- (12,2.8);
    \draw[thick] (0,2.6) -- (12,2.6);
    \node at (0.5,2.3) {$\vdots$};
    
    \draw[thick] (0,1) -- (12,1);
    \draw[thick] (0,0.8) -- (12,0.8);
    \draw[thick] (0,0.6) -- (12,0.6);
    \node at (0.5,0.3) {$\vdots$};
    
    % Gates
    \node[draw, thick, minimum width=1.5cm, minimum height=2.5cm, fill=blue!20] (h) at (2,2) {$H^{\otimes n}$};
    \node[draw, thick, minimum width=3cm, minimum height=2.5cm, fill=red!20] (mod) at (6,2) {
      \begin{tabular}{c}
        Modular\\Exponentiation\\$U_a|x\rangle|0\rangle = |x\rangle|a^x \bmod n\rangle$
      \end{tabular}
    };
    \node[draw, thick, minimum width=1.5cm, minimum height=1.2cm, fill=green!20] (qft) at (10,3) {QFT$^{-1}$};
    
    % Measurement
    \draw[thick] (11.5,3) -- (11.5,2.5) -- (12,2.5) -- (12,3.5) -- (11.5,3.5) -- (11.5,3);
    \node[right] at (12.2,3) {Measure};
    
    % Annotations
    \node[below, blue] at (2,0) {Superposition};
    \node[below, red] at (6,-0.5) {Main computational cost};
    \node[below, green] at (10,0) {Extract period};
  \end{tikzpicture}
  
  \vspace{0.5cm}
  
  \begin{greenBox}
    \textbf{Resource Requirements:}
    \begin{itemize}
      \item $\sim 3n$ logical qubits for $n$-bit RSA modulus
      \item $\sim n^3$ gates (mostly in modular exponentiation)
      \item Polynomial depth $O(n^3)$
    \end{itemize}
  \end{greenBox}
  
  \teaching{Show it's a real, implementable algorithm}
\end{frame}

% Slide 2.1.7: The Devastating Efficiency
\begin{frame}
  \slideheader{The Devastating Efficiency}
  
  \begin{center}
    \textbf{\large Classical vs Quantum Time Complexity}
  \end{center}
  
  \vspace{0.3cm}
  
  \begin{table}
    \centering
    \begin{tabular}{|l|r|r|c|}
      \hline
      \textbf{RSA Size} & \textbf{Classical Time} & \textbf{Quantum Time} & \textbf{Speedup} \\
      \hline
      RSA-1024 & $10^6$ years & 3 hours & $10^{10}\times$ \\
      RSA-2048 & $10^{20}$ years & 8 hours & $10^{24}\times$ \\
      RSA-4096 & $10^{40}$ years & 15 hours & $10^{44}\times$ \\
      RSA-8192 & $10^{80}$ years & 30 hours & $10^{84}\times$ \\
      \hline
    \end{tabular}
  \end{table}
  
  \vspace{0.5cm}
  
  \begin{redBox}
    \textbf{No escape through bigger keys!}
    \begin{itemize}
      \item Classical: Exponential in key size
      \item Quantum: Polynomial in key size
      \item Doubling key size: Classical $\times 10^{20}$, Quantum $\times 2$
    \end{itemize}
  \end{redBox}
  
  \vspace{0.3cm}
  
  \begin{tikzpicture}[scale=0.6]
    \draw[thick, ->] (0,0) -- (8,0) node[right] {Key size};
    \draw[thick, ->] (0,0) -- (0,5) node[above] {Time};
    
    % Classical exponential curve
    \draw[red, ultra thick, domain=0:6] plot (\x, {0.1*exp(0.5*\x)});
    \node[red, right] at (6,4) {Classical};
    
    % Quantum polynomial curve
    \draw[blue, ultra thick, domain=0:7] plot (\x, {0.3*\x});
    \node[blue, right] at (7,2.1) {Quantum};
  \end{tikzpicture}
  
  \teaching{Drive home the exponential vs polynomial gap}
\end{frame}

% Slide 2.1.8: Current Quantum Hardware Reality Check
\begin{frame}
  \slideheader{Current Quantum Hardware Reality Check}
  
  \begin{columns}[T]
    \begin{column}{0.55\textwidth}

    \textbf{Requirements for Breaking RSA-2048:}
    
    \begin{itemize}
      \item \textbf{Logical qubits:} $\sim$4,000
      \item \textbf{Error rate:} $< 10^{-10}$
      \item \textbf{Gate depth:} $\sim 10^9$
    \end{itemize}
    
    \vspace{0.3cm}
    
    \textbf{Physical Reality:}
    
    \begin{itemize}
      \item Logical $\rightarrow$ Physical: $\sim$10,000:1
      \item \textbf{Total needed:} 20-40 million physical qubits
      \item \textbf{Current record:} $\sim$1,000 qubits
      \item \textbf{Gap:} 4-5 orders of magnitude
    \end{itemize}
    
    \vspace{0.3cm}
    
    \begin{blueBox}
      \textbf{Timeline Estimate:}\\
      10-20 years for cryptographically relevant quantum computers
    \end{blueBox}
  
    \end{column}
    \begin{column}{0.45\textwidth}

    \begin{tikzpicture}[scale=0.7]
      % Log scale plot
      \draw[thick, ->] (0,0) -- (5,0) node[right] {Year};
      \draw[thick, ->] (0,0) -- (0,5) node[above] {Qubits (log)};
      
      % Y-axis labels
      \node[left] at (0,1) {$10^2$};
      \node[left] at (0,2) {$10^3$};
      \node[left] at (0,3) {$10^4$};
      \node[left] at (0,4) {$10^5$};
      
      % X-axis labels
      \node[below] at (1,0) {2020};
      \node[below] at (2.5,0) {2025};
      \node[below] at (4,0) {2030};
      
      % Current progress
      \draw[blue, ultra thick] plot[smooth] coordinates {
        (0,0.7) (0.5,1) (1,1.3) (1.5,1.7) (2,2) (2.5,2.3)
      };
      
      % Projected (dashed)
      \draw[blue, ultra thick, dashed] plot[smooth] coordinates {
        (2.5,2.3) (3,2.7) (3.5,3.3) (4,4)
      };
      
      % Danger threshold
      \draw[red, ultra thick, dashed] (0,4.5) -- (5,4.5);
      \node[red, right] at (5,4.5) {Crypto-relevant};
      
      \node[green] at (2.5,1) {\textbf{We are here}};
    \end{tikzpicture}
  
    \end{column}
  \end{columns}
  
  \teaching{Balanced view - real threat but time to prepare}
\end{frame}

% Slide 2.1.9: Breaking Signatures Too!
\begin{frame}
  \slideheader{Breaking Signatures Too!}
  
  \begin{center}
    \textbf{\large Same Quantum Algorithms Break All Classical Public Key Crypto}
  \end{center}
  
  \vspace{0.5cm}
  
  \begin{table}
    \centering
    \begin{tabular}{|l|l|l|}
      \hline
      \textbf{Signature Scheme} & \textbf{Underlying Problem} & \textbf{Quantum Algorithm} \\
      \hline
      RSA Signatures & Integer Factorization & Shor's Algorithm \\
      DSA & Discrete Logarithm & Modified Shor \\
      ECDSA & Elliptic Curve DLog & Modified Shor \\
      EdDSA & Elliptic Curve DLog & Modified Shor \\
      \hline
    \end{tabular}
  \end{table}
  
  \vspace{0.5cm}
  
  \begin{redBox}
    \textbf{Impact of Signature Breaking:}
    \begin{itemize}
      \item Forge ANY digital signature
      \item Impersonate ANY website (fake certificates)
      \item Falsify software updates
      \item Compromise blockchain transactions
      \item Destroy digital trust infrastructure
    \end{itemize}
  \end{redBox}
  
  \vspace{0.3cm}
  
  \begin{greenBox}
    \textbf{Timeline:} Same as encryption - when quantum computers can break RSA-2048, they can forge RSA-2048 signatures
  \end{greenBox}
  
  \teaching{Emphasize complete failure of classical public key crypto}
\end{frame}

% Slide 2.1.10: Grover's Algorithm - The Symmetric Threat
\begin{frame}
  \slideheader{Grover's Algorithm - The Symmetric Threat}
  
  \begin{columns}[T]
    \begin{column}{0.5\textwidth}

    \textbf{Different Quantum Approach:}
    \begin{itemize}
      \item Amplitude amplification
      \item Search $N$ items in $\sqrt{N}$ time
      \item Quadratic speedup (not exponential)
    \end{itemize}
    
    \vspace{0.3cm}
    
    \textbf{Impact on Symmetric Crypto:}
    
    \begin{tabular}{|l|c|c|}
      \hline
      \textbf{Algorithm} & \textbf{Classical} & \textbf{Quantum} \\
      \hline
      AES-128 & 128-bit & 64-bit \\
      AES-256 & 256-bit & 128-bit \\
      SHA-256 & 128-bit coll. & 85-bit coll. \\
      SHA-512 & 256-bit coll. & 170-bit coll. \\
      \hline
    \end{tabular}
    
    \vspace{0.3cm}
    
    \begin{greenBox}
      \textbf{Solution:} Double key sizes\\
      AES-256 remains secure\\
      SHA-384/512 remain secure
    \end{greenBox}
  
    \end{column}
    \begin{column}{0.5\textwidth}

    \begin{tikzpicture}[scale=0.7]
      % Search space visualization
      \draw[thick] (0,0) rectangle (4,4);
      
      % Grid of items
      \foreach \x in {0.5,1.5,2.5,3.5} {
        \foreach \y in {0.5,1.5,2.5,3.5} {
          \node[circle, draw, minimum size=0.3cm] at (\x,\y) {};
        }
      }
      
      % Marked item
      \node[circle, fill=red, minimum size=0.3cm] at (2.5,1.5) {};
      
      % Classical search path
      \draw[->, blue, thick] (0.5,3.5) -- (1.5,3.5) -- (2.5,3.5) -- (3.5,3.5) -- (3.5,2.5) -- (2.5,2.5) -- (1.5,2.5) -- (0.5,2.5) -- (0.5,1.5) -- (1.5,1.5) -- (2.5,1.5);
      \node[blue, below] at (2,-0.5) {Classical: $O(N)$ steps};
      
      % Quantum amplitude
      \begin{scope}[shift={(6,0)}]
        \draw[thick, ->] (0,0) -- (4,0) node[right] {Item};
        \draw[thick, ->] (0,0) -- (0,3) node[above] {Amplitude};
        
        % Initial uniform
        \draw[gray, dashed] (0,0.5) -- (4,0.5);
        
        % After Grover iterations
        \draw[red, ultra thick] (2.5,2.5) -- (2.5,0);
        
        \node[below] at (2,-0.5) {Quantum: $O(\sqrt{N})$ steps};
      \end{scope}
    \end{tikzpicture}
  
    \end{column}
  \end{columns}
  
  \teaching{Complete the quantum algorithm picture}
\end{frame}

% Slide 2.1.11: Why These Quantum Advantages?
\begin{frame}
  \slideheader{Why These Quantum Advantages?}
  
  \begin{center}
    \textbf{\large What Makes Some Problems Quantum-Easy?}
  \end{center}
  
  \vspace{0.5cm}
  
  \begin{columns}[T]
    \begin{column}{0.5\textwidth}

    \begin{blueBox}
      \textbf{Shor Exploits:}
      \begin{itemize}
        \item Hidden periodic structure
        \item Fourier analysis natural for QC
        \item Interference extracts global property
      \end{itemize}
    \end{blueBox}
    
    \vspace{0.3cm}
    
    \begin{blueBox}
      \textbf{Grover Exploits:}
      \begin{itemize}
        \item Superposition over search space
        \item Amplitude amplification
        \item Quadratic speedup only
      \end{itemize}
    \end{blueBox}
  
    \end{column}
    \begin{column}{0.5\textwidth}

    \begin{tikzpicture}[scale=0.6]
      % Venn diagram
      \draw[thick, fill=blue!20] (0,0) circle (2cm);
      \draw[thick, fill=red!20] (2,0) circle (2cm);
      
      \node at (-1,0) {\small \textbf{Periodic}};
      \node at (3,0) {\small \textbf{Search}};
      \node at (1,0) {\small \textbf{Both}};
      
      % Examples
      \node[below] at (-1,-2.5) {\tiny Factoring, DLog};
      \node[below] at (3,-2.5) {\tiny Database search};
    \end{tikzpicture}
  
    \end{column}
  \end{columns}
  
  \vspace{0.5cm}
  
  \begin{redBox}
    \textbf{Key Insight:} Not all problems have quantum speedup!
    \begin{itemize}
      \item Need specific structure (periodicity, search)
      \item Most problems: No quantum advantage
      \item This gives us hope for post-quantum crypto!
    \end{itemize}
  \end{redBox}
  
  \teaching{Transition to quantum-resistant problems}
\end{frame}

% ==================================================
% Section 2.2: Enter the Lattice
% ==================================================

\section{Enter the Lattice}

% Slide 2.2.1: From Quantum Vulnerability to Quantum Resistance
\begin{frame}
  \slideheader{From Quantum Vulnerability to Quantum Resistance}
  
  \begin{center}
    \textbf{\large The Search for Quantum-Hard Problems}
  \end{center}
  
  \vspace{0.5cm}
  
  \begin{blueBox}
    \textbf{Requirements for Post-Quantum Security:}
    \begin{itemize}
      \item[$\times$] No hidden periodicity (Shor doesn't apply)
      \item[$\times$] No efficient quantum algorithm known
      \item[\checkmark] Still usable for cryptography
      \item[\checkmark] Efficient classical algorithms for legitimate users
      \item[\checkmark] Decades of cryptanalysis attempts
    \end{itemize}
  \end{blueBox}
  
  \vspace{0.5cm}
  
  \pause
  
  \begin{center}
    \begin{tikzpicture}[scale=1.2]
      % Arrow pointing down
      \draw[->, ultra thick, green!60!black] (0,0) -- (0,-1);
      
      % Text
      \node[font=\Large\bfseries] at (0,-1.5) {Enter: LATTICE PROBLEMS};
      
      % Lattice visualization
      \begin{scope}[shift={(0,-3)}]
        \foreach \x in {-2,-1,0,1,2} {
          \foreach \y in {-1,0,1} {
            \node[circle, fill=black, inner sep=1.5pt] at (\x*0.6,\y*0.6) {};
          }
        }
        % Connect some points
        \draw[thick, blue] (-1.2,-0.6) -- (-0.6,0) -- (0,0.6) -- (0.6,0) -- (1.2,-0.6);
        \draw[thick, blue] (-1.2,0) -- (-0.6,0.6) -- (0,0) -- (0.6,-0.6) -- (1.2,0);
      \end{scope}
    \end{tikzpicture}
  \end{center}
  
  \teaching{Introduce lattices as the solution}
\end{frame}

% Slide 2.2.2: Lattices - Your Physics Background Helps!
\begin{frame}
  \slideheader{Lattices - Your Physics Background Helps!}
  
  \begin{columns}[T]
    \begin{column}{0.5\textwidth}

    \textbf{Mathematical Definition:}
    \begin{align}
      \mathcal{L} = \left\{ \sum_{i=1}^n a_i \mathbf{b}_i : a_i \in \mathbb{Z} \right\}
    \end{align}
    
    \vspace{0.3cm}
    
    \textbf{You Already Know This From:}
    \begin{itemize}
      \item Crystal structures (Bravais lattices)
      \item Reciprocal lattices
      \item Brillouin zones
      \item Miller indices
      \item X-ray diffraction
    \end{itemize}
    
    \vspace{0.3cm}
    
    \begin{greenBox}
      \textbf{Your Advantage:}\\
      Geometric intuition from condensed matter physics!
    \end{greenBox}
  
    \end{column}
    \begin{column}{0.5\textwidth}

    \begin{tikzpicture}[scale=1.2]
      % 2D lattice
      \draw[thick, ->] (-0.5,0) -- (3.5,0) node[right] {$\mathbf{b}_1$};
      \draw[thick, ->] (0,-0.5) -- (0,3) node[above] {$\mathbf{b}_2$};
      
      % Basis vectors
      \draw[ultra thick, blue, ->] (0,0) -- (1,0.2) node[midway, below] {$\mathbf{b}_1$};
      \draw[ultra thick, blue, ->] (0,0) -- (0.3,1) node[midway, left] {$\mathbf{b}_2$};
      
      % Lattice points
      \foreach \x in {-1,0,1,2,3} {
        \foreach \y in {-1,0,1,2} {
          \pgfmathsetmacro{\px}{\x*1 + \y*0.3}
          \pgfmathsetmacro{\py}{\x*0.2 + \y*1}
          \ifthenelse{\lengthtest{\px pt < 3.5pt} \AND \lengthtest{\py pt < 3pt} \AND \lengthtest{\px pt > -0.5pt} \AND \lengthtest{\py pt > -0.5pt}}{
            \node[circle, fill=black, inner sep=1.5pt] at (\px,\py) {};
          }{}
        }
      }
      
      % Example lattice vector
      \draw[red, ultra thick, ->] (0,0) -- (2.3,2.2);
      \node[red, right] at (2.3,2.2) {$2\mathbf{b}_1 + 2\mathbf{b}_2$};
    \end{tikzpicture}
    
    \vspace{0.3cm}
    
    \textbf{Crystal Structure Analogy:}\\
    FCC, BCC, hexagonal - all lattices!
  
    \end{column}
  \end{columns}
  
  \teaching{Connect to familiar physics concepts}
\end{frame}

% Slide 2.2.3: Good Basis vs Bad Basis
\begin{frame}
  \slideheader{Good Basis vs Bad Basis}
  
  \begin{center}
    \textbf{\large Same Lattice, Different Bases}
  \end{center}
  
  \vspace{0.3cm}
  
  \begin{columns}[T]
    \begin{column}{0.5\textwidth}

    \begin{center}
      \textbf{Good Basis}
    \end{center}
    
    \begin{tikzpicture}[scale=1.5]
      % Good basis vectors
      \draw[ultra thick, blue, ->] (0,0) -- (1,0) node[midway, below] {$\mathbf{b}_1$};
      \draw[ultra thick, blue, ->] (0,0) -- (0.1,0.9) node[midway, left] {$\mathbf{b}_2$};
      
      % Lattice points
      \foreach \x in {-1,0,1,2} {
        \foreach \y in {-1,0,1,2} {
          \pgfmathsetmacro{\px}{\x + 0.1*\y}
          \pgfmathsetmacro{\py}{0.9*\y}
          \node[circle, fill=gray, inner sep=1pt] at (\px,\py) {};
        }
      }
      
      % Properties
      \node[below] at (1,-0.5) {
        \begin{tabular}{l}
          Nearly orthogonal\\
          Short vectors\\
          Easy to work with
        \end{tabular}
      };
    \end{tikzpicture}
  
    \end{column}
    \begin{column}{0.5\textwidth}

    \begin{center}
      \textbf{Bad Basis}
    \end{center}
    
    \begin{tikzpicture}[scale=1.5]
      % Bad basis vectors
      \draw[ultra thick, red, ->] (0,0) -- (2.1,0.9) node[midway, below] {$\mathbf{b}_1'$};
      \draw[ultra thick, red, ->] (0,0) -- (1.9,1.8) node[midway, above] {$\mathbf{b}_2'$};
      
      % Same lattice points
      \foreach \x in {-1,0,1,2} {
        \foreach \y in {-1,0,1,2} {
          \pgfmathsetmacro{\px}{\x + 0.1*\y}
          \pgfmathsetmacro{\py}{0.9*\y}
          \node[circle, fill=gray, inner sep=1pt] at (\px,\py) {};
        }
      }
      
      % Properties
      \node[below] at (1,-0.5) {
        \begin{tabular}{l}
          Highly skewed\\
          Long vectors\\
          Hard to analyze
        \end{tabular}
      };
    \end{tikzpicture}
  
    \end{column}
  \end{columns}
  
  \vspace{0.5cm}
  
  \begin{redBox}
    \textbf{Key Computational Problem:}\\
    Given bad basis → Find good basis = HARD!
  \end{redBox}
  
  \begin{greenBox}
    \textbf{Crystallography Analogy:}\\
    Primitive cell vs Conventional cell - finding the ``natural'' basis is non-trivial
  \end{greenBox}
  
  \teaching{Introduce the basis reduction challenge}
\end{frame}

% Slide 2.2.4: The Dimension Explosion
\begin{frame}
  \slideheader{The Dimension Explosion}
  
  \begin{columns}[T]
    \begin{column}{0.5\textwidth}

    \textbf{Your 3D Intuition:}
    \begin{itemize}
      \item Vectors have clear ``direction''
      \item Can visualize shortest vector
      \item Volume is... volume
      \item Few ``nearly shortest'' vectors
    \end{itemize}
    
    \vspace{0.5cm}
    
    \textbf{In 500+ Dimensions:}
    \begin{itemize}
      \item Volume concentrates at surface
      \item All vectors nearly orthogonal
      \item Exponentially many short vectors
      \item Geometry becomes counterintuitive
    \end{itemize}
  
    \end{column}
    \begin{column}{0.5\textwidth}

    \begin{tikzpicture}[scale=0.8]
      % Sphere volume concentration
      \draw[thick] (0,0) circle (2cm);
      \draw[thick, fill=blue!30] (0,0) circle (1.8cm);
      \draw[thick, fill=blue!60] (0,0) circle (1.6cm);
      
      \node[below] at (0,-2.5) {$n = 3$};
      \node at (0,0) {90\% volume};
      
      \begin{scope}[shift={(5,0)}]
        \draw[thick] (0,0) circle (2cm);
        \draw[thick, fill=blue!90] (0,0) circle (1.99cm);
        
        \node[below] at (0,-2.5) {$n = 500$};
        \node at (0,0) {\tiny 99.9\% at surface!};
      \end{scope}
    \end{tikzpicture}
    
    \vspace{0.5cm}
    
    \begin{redBox}
      \textbf{Curse of Dimensionality:}\\
      Algorithms that work in 3D fail spectacularly in 500D
    \end{redBox}
  
    \end{column}
  \end{columns}
  
  \teaching{Explain why high dimensions matter}
\end{frame}

% Slide 2.2.5: The Shortest Vector Problem (SVP)
\begin{frame}
  \slideheader{The Shortest Vector Problem (SVP)}
  
  \begin{center}
    \textbf{\large Core Lattice Problem \#1}
  \end{center}
  
  \vspace{0.3cm}
  
  \begin{blueBox}
    \textbf{Given:} A basis $B = \{\mathbf{b}_1, \ldots, \mathbf{b}_n\}$ of lattice $\mathcal{L}$\\
    \textbf{Find:} The shortest non-zero vector $\mathbf{v} \in \mathcal{L}$
  \end{blueBox}
  
  \vspace{0.3cm}
  
  \begin{columns}[T]
    \begin{column}{0.5\textwidth}

    \begin{tikzpicture}[scale=1.2]
      % Bad basis
      \draw[thick, ->] (0,0) -- (2.5,0.3) node[right] {$\mathbf{b}_1$};
      \draw[thick, ->] (0,0) -- (2.2,1.8) node[above] {$\mathbf{b}_2$};
      
      % Lattice points
      \foreach \x in {-1,0,1,2} {
        \foreach \y in {-1,0,1,2} {
          \pgfmathsetmacro{\px}{\x*2.5 + \y*2.2}
          \pgfmathsetmacro{\py}{\x*0.3 + \y*1.8}
          \ifthenelse{\lengthtest{\px pt < 5pt} \AND \lengthtest{\py pt < 4pt} \AND \lengthtest{\px pt > -1pt} \AND \lengthtest{\py pt > -1pt}}{
            \node[circle, fill=gray, inner sep=1.5pt] at (\px,\py) {};
          }{}
        }
      }
      
      % Shortest vector (not obvious!)
      \draw[ultra thick, red, ->] (0,0) -- (0.3,1.5);
      \node[red, left] at (0.3,1.5) {$\mathbf{v}_{min}$};
      
      % Circle showing minimum distance
      \draw[dashed, red] (0,0) circle (1.52cm);
    \end{tikzpicture}
  
    \end{column}
    \begin{column}{0.5\textwidth}

    \textbf{Complexity:}
    \begin{itemize}
      \item NP-hard to solve exactly
      \item NP-hard to approximate within factor $n^c$
      \item Best algorithms: $2^{O(n)}$ time
    \end{itemize}
    
    \vspace{0.3cm}
    
    \textbf{Why it's hard:}
    \begin{itemize}
      \item Can't just try basis vectors
      \item Need integer combinations
      \item Exponentially many to check
    \end{itemize}
  
    \end{column}
  \end{columns}
  
  \teaching{Define the core computational problem}
\end{frame}

% Slide 2.2.6: The Closest Vector Problem (CVP)
\begin{frame}
  \slideheader{The Closest Vector Problem (CVP)}
  
  \begin{center}
    \textbf{\large Core Lattice Problem \#2}
  \end{center}
  
  \vspace{0.3cm}
  
  \begin{blueBox}
    \textbf{Given:} Lattice $\mathcal{L}$ and target point $\mathbf{t} \notin \mathcal{L}$\\
    \textbf{Find:} Closest lattice point $\mathbf{v} \in \mathcal{L}$ to $\mathbf{t}$
  \end{blueBox}
  
  \vspace{0.3cm}
  
  \begin{columns}[T]
    \begin{column}{0.5\textwidth}

    \begin{tikzpicture}[scale=1.2]
      % Lattice points
      \foreach \x in {-1,0,1,2,3} {
        \foreach \y in {-1,0,1,2} {
          \node[circle, fill=gray, inner sep=1.5pt] at (\x,\y) {};
        }
      }
      
      % Target point
      \node[star, star points=5, star point ratio=0.5, fill=blue, minimum size=0.3cm] (target) at (1.7,1.3) {};
      \node[blue, above] at (target) {$\mathbf{t}$};
      
      % Closest lattice point
      \node[circle, fill=red, inner sep=2pt] at (2,1) {};
      \draw[red, thick, ->] (1.7,1.3) -- (2,1);
      
      % Voronoi cell indication
      \draw[dashed, green] (1.5,0.5) -- (2.5,0.5) -- (2.5,1.5) -- (1.5,1.5) -- cycle;
    \end{tikzpicture}
  
    \end{column}
    \begin{column}{0.5\textwidth}

    \textbf{Connection to Coding:}
    \begin{itemize}
      \item Lattice = Error correcting code
      \item Target = Received noisy codeword
      \item CVP = Decode to nearest codeword
    \end{itemize}
    
    \vspace{0.3cm}
    
    \textbf{Also NP-hard!}
    \begin{itemize}
      \item Even harder than SVP
      \item Used directly in crypto
    \end{itemize}
  
    \end{column}
  \end{columns}
  
  \vspace{0.3cm}
  
  \begin{greenBox}
    \textbf{Physical Analogy:}\\
    Finding which unit cell a perturbed atom belongs to
  \end{greenBox}
  
  \teaching{Connect to error correction}
\end{frame}

% Slide 2.2.7: Classical Algorithms for Lattice Problems
\begin{frame}
  \slideheader{Classical Algorithms for Lattice Problems}
  
  \begin{center}
    \textbf{\large The Algorithm Zoo}
  \end{center}
  
  \vspace{0.3cm}
  
  \begin{table}
    \centering
    \begin{tabular}{|l|l|l|l|}
      \hline
      \textbf{Algorithm} & \textbf{Year} & \textbf{Time} & \textbf{Quality} \\
      \hline
      LLL & 1982 & Polynomial & Exponential approx \\
      BKZ-$k$ & 1994 & $2^{O(k)}$ & $(2k)^{n/2k}$ approx \\
      Sieving & 2008 & $2^{0.292n + o(n)}$ & Exact SVP \\
      Enumeration & 2010 & $2^{O(n \log n)}$ & Exact SVP \\
      \hline
    \end{tabular}
  \end{table}
  
  \vspace{0.5cm}
  
  \begin{itemize}
    \item \textbf{LLL:} Fast but poor quality (sufficient for some attacks)
    \item \textbf{BKZ:} Time-quality tradeoff via blocksize $k$
    \item \textbf{Sieving:} Best asymptotic complexity
    \item \textbf{Key point:} All exponential for good approximations!
  \end{itemize}
  
  \vspace{0.3cm}
  
  \begin{tikzpicture}[scale=0.7]
    \draw[thick, ->] (0,0) -- (6,0) node[right] {Dimension $n$};
    \draw[thick, ->] (0,0) -- (0,4) node[above] {Time (log scale)};
    
    % LLL - polynomial
    \draw[green, ultra thick] plot[smooth] coordinates {(0,0.5) (2,0.7) (4,0.9) (6,1.1)};
    \node[green, right] at (6,1.1) {LLL};
    
    % BKZ - exponential
    \draw[blue, ultra thick] plot[smooth] coordinates {(0,0.5) (2,1.5) (4,2.5) (6,3.5)};
    \node[blue, right] at (6,3.5) {BKZ};
    
    % Sieving - exponential but better
    \draw[red, ultra thick] plot[smooth] coordinates {(0,0.5) (2,1.2) (4,2) (6,2.8)};
    \node[red, right] at (6,2.8) {Sieving};
  \end{tikzpicture}
  
  \teaching{Establish classical hardness baseline}
\end{frame}

% Slide 2.2.8: Quantum (Non-)Algorithms for Lattices
\begin{frame}
  \slideheader{Quantum (Non-)Algorithms for Lattices}
  
  \begin{center}
    \textbf{\large The Quantum Speedup That Wasn't}
  \end{center}
  
  \vspace{0.5cm}
  
  \begin{table}
    \centering
    \begin{tabular}{|l|c|c|c|}
      \hline
      \textbf{Algorithm} & \textbf{Classical} & \textbf{Quantum} & \textbf{Speedup} \\
      \hline
      Best SVP solver & $2^{0.292n}$ & $2^{0.265n}$ & $\approx 1.019^n$ \\
      \hline
      Compare: Factoring & $2^{n^{1/3}}$ & $\text{poly}(n)$ & Exponential! \\
      \hline
    \end{tabular}
  \end{table}
  
  \vspace{0.5cm}
  
  \begin{columns}[T]
    \begin{column}{0.5\textwidth}

    \begin{redBox}
      \textbf{Tiny Quantum Advantage:}
      \begin{itemize}
        \item Classical: $2^{0.292n}$
        \item Quantum: $2^{0.265n}$
        \item Speedup: $2^{0.027n} \approx 1.019^n$
        \item For $n=500$: $1.019^{500} \approx 13,000\times$
      \end{itemize}
    \end{redBox}
    
    \vspace{0.3cm}
    
    Compare to Shor:
    \begin{itemize}
      \item Classical: $2^{n^{1/3}}$
      \item Quantum: $n^3$
      \item Speedup: $2^{n^{1/3}}/n^3$ (astronomical!)
    \end{itemize}
  
    \end{column}
    \begin{column}{0.5\textwidth}

    \begin{tikzpicture}[scale=0.7]
      % Comparison chart
      \draw[thick, ->] (0,0) -- (5,0) node[right] {Security param};
      \draw[thick, ->] (0,0) -- (0,4) node[above] {Speedup};
      
      % Lattice speedup (basically flat)
      \draw[blue, ultra thick] (0,0.2) -- (5,0.4);
      \node[blue, above] at (2.5,0.4) {Lattice: $1.019^n$};
      
      % Factoring speedup (exponential)
      \draw[red, ultra thick, domain=0.5:4.5] plot (\x, {0.3*exp(0.5*\x)});
      \node[red, right] at (4.5,3.5) {Factoring};
    \end{tikzpicture}
  
    \end{column}
  \end{columns}
  
  \vspace{0.3cm}
  
  \begin{greenBox}
    \textbf{Bottom Line:} No meaningful quantum speedup for lattices!
  \end{greenBox}
  
  \teaching{Key message - lattices resist quantum attack}
\end{frame}

% Slide 2.2.9: Why No Quantum Breakthrough?
\begin{frame}
  \slideheader{Why No Quantum Breakthrough?}
  
  \begin{center}
    \textbf{\large What Makes Lattices Quantum-Resistant?}
  \end{center}
  
  \vspace{0.5cm}
  
  \begin{columns}[T]
    \begin{column}{0.5\textwidth}

    \textbf{No Hidden Periodicity:}
    \begin{itemize}
      \item Shor needs periodic structure
      \item Lattice problems are geometric
      \item No algebraic pattern to exploit
    \end{itemize}
    
    \vspace{0.3cm}
    
    \textbf{No Algebraic Structure:}
    \begin{itemize}
      \item Can't use quantum Fourier transform
      \item No group structure to leverage
      \item Pure geometric optimization
    \end{itemize}
    
    \vspace{0.3cm}
    
    \textbf{High-Dimensional Geometry:}
    \begin{itemize}
      \item Quantum walks don't help much
      \item No efficient amplitude amplification
      \item Curse of dimensionality affects QC too!
    \end{itemize}
  
    \end{column}
    \begin{column}{0.5\textwidth}

    \begin{tikzpicture}[scale=0.8]
      % Timeline of attempts
      \draw[thick, ->] (0,0) -- (0,5) node[above] {Years};
      
      \node[left] at (0,0.5) {2008};
      \node[left] at (0,1.5) {2011};
      \node[left] at (0,2.5) {2015};
      \node[left] at (0,3.5) {2020};
      \node[left] at (0,4.5) {2024};
      
      % Attempts
      \node[right, red] at (0.2,0.5) {$\times$ Quantum sieve?};
      \node[right, red] at (0.2,1.5) {$\times$ Quantum enumeration?};
      \node[right, red] at (0.2,2.5) {$\times$ Novel approaches?};
      \node[right, red] at (0.2,3.5) {$\times$ Hybrid algorithms?};
      \node[right, red] at (0.2,4.5) {Still no breakthrough!};
    \end{tikzpicture}
  
    \end{column}
  \end{columns}
  
  \vspace{0.3cm}
  
  \begin{blueBox}
    \textbf{15+ Years of Research:}\\
    Despite intense effort, no exponential quantum speedup found for lattice problems
  \end{blueBox}
  
  \teaching{Build confidence in quantum resistance}
\end{frame}

% Slide 2.2.10: Concrete Security Numbers
\begin{frame}
  \slideheader{Concrete Security Numbers}
  
  \begin{center}
    \textbf{\large Dimension vs Security Level}
  \end{center}
  
  \vspace{0.3cm}
  
  \begin{table}
    \centering
    \begin{tabular}{|c|c|c|c|}
      \hline
      \textbf{Dimension} & \textbf{Classical} & \textbf{Quantum} & \textbf{Equivalent} \\
      \hline
      $n = 300$ & $2^{88}$ & $2^{80}$ & - \\
      $n = 500$ & $2^{146}$ & $2^{133}$ & AES-128 \\
      $n = 700$ & $2^{204}$ & $2^{186}$ & RSA-2048 \\
      $n = 1000$ & $2^{292}$ & $2^{265}$ & AES-256 \\
      \hline
    \end{tabular}
  \end{table}
  
  \vspace{0.5cm}
  
  \begin{columns}[T]
    \begin{column}{0.5\textwidth}

    \begin{greenBox}
      \textbf{Key Takeaways:}
      \begin{itemize}
        \item Still exponentially hard for QC
        \item Modest dimension increase compensates
        \item Well-understood security
      \end{itemize}
    \end{greenBox}
    
    \vspace{0.3cm}
    
    \textbf{For RSA-2048 equivalent:}
    \begin{itemize}
      \item Use lattice dimension $\approx 700$
      \item Security margin included
      \item Future-proof for decades
    \end{itemize}
  
    \end{column}
    \begin{column}{0.5\textwidth}

    \begin{tikzpicture}[scale=0.7]
      % Security level comparison
      \draw[thick, ->] (0,0) -- (5,0) node[right] {Dimension};
      \draw[thick, ->] (0,0) -- (0,4) node[above] {Security (bits)};
      
      % Classical security
      \draw[blue, ultra thick] plot[smooth] coordinates {
        (0,0) (1,0.88) (2,1.46) (3,2.04) (4,2.92)
      };
      \node[blue, above] at (4,2.92) {Classical};
      
      % Quantum security
      \draw[red, ultra thick] plot[smooth] coordinates {
        (0,0) (1,0.8) (2,1.33) (3,1.86) (4,2.65)
      };
      \node[red, below] at (4,2.65) {Quantum};
      
      % Markers
      \draw[dashed] (2,0) -- (2,1.46);
      \node[below] at (2,0) {500};
      \draw[dashed] (3,0) -- (3,2.04);
      \node[below] at (3,0) {700};
    \end{tikzpicture}
  
    \end{column}
  \end{columns}
  
  \teaching{Quantify the security with concrete numbers}
\end{frame}

% Slide 2.2.11: The Lattice Advantage Summary
\begin{frame}
  \slideheader{The Lattice Advantage Summary}
  
  \begin{center}
    \textbf{\large Why Lattices Win for Post-Quantum Crypto}
  \end{center}
  
  \vspace{0.5cm}
  
  \begin{columns}
    \begin{column}{0.5\textwidth}
      \begin{greenBox}
        \textbf{Security Advantages:}
        \begin{itemize}
          \item[\checkmark] No quantum algorithm breakthrough
          \item[\checkmark] 15+ years of quantum cryptanalysis
          \item[\checkmark] Worst-case hardness guarantees
          \item[\checkmark] Well-understood security
        \end{itemize}
      \end{greenBox}
      
      \vspace{0.3cm}
      
      \begin{greenBox}
        \textbf{Practical Advantages:}
        \begin{itemize}
          \item[\checkmark] Efficient implementations
          \item[\checkmark] Multiple crypto primitives
          \item[\checkmark] Parallelizable operations
          \item[\checkmark] Hardware-friendly
        \end{itemize}
      \end{greenBox}
    \end{column}
    
    \begin{column}{0.5\textwidth}
      \begin{tikzpicture}[scale=0.8]
        % Trophy icon
        \draw[thick, fill=yellow!50] (-1,2) -- (-1.2,1) -- (-0.5,1) -- (-0.5,0) -- (0.5,0) -- (0.5,1) -- (1.2,1) -- (1,2) -- cycle;
        \draw[thick, fill=yellow!50] (-1,2) .. controls (-1,2.5) and (1,2.5) .. (1,2);
        
        % Handles
        \draw[thick] (-1,1.7) .. controls (-1.5,1.7) and (-1.5,1.3) .. (-1.2,1.3);
        \draw[thick] (1,1.7) .. controls (1.5,1.7) and (1.5,1.3) .. (1.2,1.3);
        
        % Number 1
        \node[font=\Large\bfseries] at (0,1.5) {1};
        
        % Text below
        \node[below, align=center, font=\large\bfseries] at (0,-0.5) {
          Lattices:\\
          The Winner of\\
          Post-Quantum\\
          Cryptography
        };
      \end{tikzpicture}
    \end{column}
  \end{columns}
  
  \vspace{0.5cm}
  
  \begin{redBox}
    \centering
    \textbf{Ready for Hour 3:} How do we actually build cryptosystems from lattices?
  \end{redBox}
  
  \teaching{Summarize lattice advantages and transition to Hour 3}
\end{frame}

% ==================================================
% END OF HOUR 2
% ==================================================

% ==================================================
% HOUR 3: SOLUTIONS AND DEPLOYMENT
% ==================================================

\section{Learning With Errors - The Bridge}

% Slide 3.1.1: From Abstract Lattices to Practical Crypto
\begin{frame}
  \slideheader{From Abstract Lattices to Practical Crypto}
  
  \begin{center}
    \textbf{\large The Challenge: How to Use Geometric Hardness?}
  \end{center}
  
  \vspace{0.5cm}
  
  \begin{tikzpicture}[scale=1.2]
    % Left: Abstract lattice
    \node[draw, thick, minimum width=3cm, minimum height=2cm, fill=blue!20] (lattice) at (-4,0) {
      \begin{tabular}{c}
        \textbf{Lattice}\\
        \textbf{Problems}\\
        (SVP, CVP)
      \end{tabular}
    };
    
    % Middle: Bridge
    \node[draw, ultra thick, minimum width=2.5cm, minimum height=2cm, fill=yellow!30] (lwe) at (0,0) {
      \begin{tabular}{c}
        \textbf{Learning}\\
        \textbf{With}\\
        \textbf{Errors}
      \end{tabular}
    };
    
    % Right: Cryptosystems
    \node[draw, thick, minimum width=3cm, minimum height=2cm, fill=green!20] (crypto) at (4,0) {
      \begin{tabular}{c}
        \textbf{Practical}\\
        \textbf{Cryptosystems}\\
        (Encryption, Sigs)
      \end{tabular}
    };
    
    % Arrows
    \draw[->, ultra thick] (lattice) -- (lwe) node[midway, above] {Reduction};
    \draw[->, ultra thick] (lwe) -- (crypto) node[midway, above] {Construction};
    
    % Labels
    \node[below] at (-4,-1.5) {\small Geometric};
    \node[below] at (0,-1.5) {\small Algebraic};
    \node[below] at (4,-1.5) {\small Practical};
  \end{tikzpicture}
  
  \vspace{0.5cm}
  
  \begin{blueBox}
    \textbf{LWE:} The crucial bridge between abstract lattice problems and usable cryptography
  \end{blueBox}
  
  \teaching{Introduce LWE as the key enabling construction}
\end{frame}

% Slide 3.1.2: The LWE Insight - Linear Algebra Plus Noise
\begin{frame}
  \slideheader{The LWE Insight - Linear Algebra Plus Noise}
  
  \begin{columns}[T]
    \begin{column}{0.5\textwidth}

    \textbf{Without Noise:}
    \begin{align}
      \mathbf{b} = A\mathbf{s}
    \end{align}
    
    \begin{itemize}
      \item Linear system of equations
      \item Solve by Gaussian elimination
      \item Polynomial time: $O(n^3)$
      \item Completely insecure!
    \end{itemize}
    
    \vspace{0.5cm}
    
    \textbf{With Small Noise:}
    \begin{align}
      \mathbf{b} = A\mathbf{s} + \mathbf{e}
    \end{align}
    
    \begin{itemize}
      \item Noisy linear system
      \item Gaussian elimination fails!
      \item Best algorithms: exponential
      \item Secure foundation
    \end{itemize}
  
    \end{column}
    \begin{column}{0.5\textwidth}

    \begin{tikzpicture}[scale=0.8]
      % Clean system
      \node at (0,3.5) {\textbf{Clean System}};
      \draw[thick] (0,2.5) -- (2,2.5) -- (2,0) -- (0,0) -- cycle;
      \node at (1,1.25) {$A$};
      
      \node at (2.5,1.25) {$\times$};
      
      \draw[thick] (3,2) -- (3.5,2) -- (3.5,0.5) -- (3,0.5) -- cycle;
      \node at (3.25,1.25) {$\mathbf{s}$};
      
      \node at (4,1.25) {$=$};
      
      \draw[thick] (4.5,2) -- (5,2) -- (5,0.5) -- (4.5,0.5) -- cycle;
      \node at (4.75,1.25) {$\mathbf{b}$};
      
      % Noisy system
      \node at (0,-0.5) {\textbf{Noisy System}};
      \draw[thick] (0,-1.5) -- (2,-1.5) -- (2,-4) -- (0,-4) -- cycle;
      \node at (1,-2.75) {$A$};
      
      \node at (2.5,-2.75) {$\times$};
      
      \draw[thick] (3,-2) -- (3.5,-2) -- (3.5,-3.5) -- (3,-3.5) -- cycle;
      \node at (3.25,-2.75) {$\mathbf{s}$};
      
      \node at (4,-2.75) {$+$};
      
      \draw[thick, fill=red!30] (4.5,-2) -- (5,-2) -- (5,-3.5) -- (4.5,-3.5) -- cycle;
      \node at (4.75,-2.75) {$\mathbf{e}$};
      
      \node at (5.5,-2.75) {$=$};
      
      \draw[thick] (6,-2) -- (6.5,-2) -- (6.5,-3.5) -- (6,-3.5) -- cycle;
      \node at (6.25,-2.75) {$\mathbf{b}$};
    \end{tikzpicture}
    
    \vspace{0.3cm}
    
    \begin{greenBox}
      \textbf{Physics Analogy:}\\
      Like thermal noise in measurements - small but ruins exact methods!
    \end{greenBox}
  
    \end{column}
  \end{columns}
  
  \teaching{Core concept - noise transforms easy to hard}
\end{frame}

% Slide 3.1.3: LWE Formal Definition
\begin{frame}
  \slideheader{LWE Formal Definition}
  
  \begin{center}
    \textbf{\large Learning With Errors Problem}
  \end{center}
  
  \vspace{0.3cm}
  
  \begin{blueBox}
    \textbf{LWE Distribution:}\\
    Given secret $\mathbf{s} \in \mathbb{Z}_q^n$ and error distribution $\chi$ over $\mathbb{Z}_q$:
    \begin{itemize}
      \item Sample $\mathbf{a}_i \leftarrow \mathbb{Z}_q^n$ uniformly
      \item Sample $e_i \leftarrow \chi$ (small error)
      \item Output $(\mathbf{a}_i, b_i = \langle \mathbf{a}_i, \mathbf{s} \rangle + e_i \bmod q)$
    \end{itemize}
  \end{blueBox}
  
  \vspace{0.3cm}
  
  \begin{redBox}
    \textbf{The LWE Problem:}\\
    Given many samples $\{(\mathbf{a}_i, b_i)\}_{i=1}^m$ from LWE distribution:\\
    $\rightarrow$ Recover the secret $\mathbf{s}$
  \end{redBox}
  
  \vspace{0.3cm}
  
  \begin{columns}
    \begin{column}{0.5\textwidth}
      \textbf{Parameters:}
      \begin{itemize}
        \item $n$: dimension (security parameter)
        \item $q$: modulus (polynomial in $n$)
        \item $\chi$: error distribution (Gaussian)
        \item $m$: number of samples
      \end{itemize}
    \end{column}
    
    \begin{column}{0.5\textwidth}
      \textbf{Hardness:}
      \begin{itemize}
        \item Best algorithms: $2^{O(n)}$ time
        \item Quantum: Still $2^{O(n)}$ time
        \item Proven as hard as lattice problems
      \end{itemize}
    \end{column}
  \end{columns}
  
  \teaching{Precise mathematical foundation}
\end{frame}

% Slide 3.1.4: Concrete LWE Example
\begin{frame}
  \slideheader{Concrete LWE Example}
  
  \begin{center}
    \textbf{\large Let's See LWE with Small Numbers}
  \end{center}
  
  Parameters: $n = 4$, $q = 97$, secret $\mathbf{s} = (3, 7, 2, 5)$
  
  \vspace{0.3cm}
  
  \begin{table}
    \centering
    \small
    \begin{tabular}{|c|c|c|c|}
      \hline
      $\mathbf{a}_i$ & $\langle \mathbf{a}_i, \mathbf{s} \rangle$ & $e_i$ & $b_i$ \\
      \hline
      $(14, 28, 91, 45)$ & $42 + 196 + 182 + 225 = 66 \bmod 97$ & $+1$ & $67$ \\
      $(73, 12, 56, 89)$ & $219 + 84 + 112 + 445 = 82 \bmod 97$ & $-1$ & $81$ \\
      $(31, 65, 22, 18)$ & $93 + 455 + 44 + 90 = 4 \bmod 97$ & $0$ & $4$ \\
      $(88, 47, 39, 71)$ & $264 + 329 + 78 + 355 = 54 \bmod 97$ & $+2$ & $56$ \\
      \hline
    \end{tabular}
  \end{table}
  
  \vspace{0.3cm}
  
  \begin{redBox}
    \textbf{Try to solve this!}
    \begin{itemize}
      \item Without errors: Easy linear algebra
      \item With errors: Which equation has $e = 0$? $\pm 1$? $\pm 2$?
      \item Can't use standard techniques!
    \end{itemize}
  \end{redBox}
  
  \vspace{0.3cm}
  
  \begin{greenBox}
    Even with tiny parameters, the problem becomes intractable quickly
  \end{greenBox}
  
  \teaching{Make it tangible with real numbers}
\end{frame}

% Slide 3.1.5: The Error Distribution
\begin{frame}
  \slideheader{The Error Distribution}
  
  \begin{center}
    \textbf{\large The Critical Component: Discrete Gaussian Noise}
  \end{center}
  
  \vspace{0.3cm}
  
  \begin{columns}[T]
    \begin{column}{0.5\textwidth}

    \textbf{Discrete Gaussian Distribution:}
    \begin{align}
      \Pr[e = x] \propto \exp\left(-\frac{x^2}{2\sigma^2}\right)
    \end{align}
    
    for $x \in \mathbb{Z}$
    
    \vspace{0.3cm}
    
    \textbf{Key Properties:}
    \begin{itemize}
      \item Center: $\mu = 0$
      \item Width: $\sigma$ (standard deviation)
      \item Support: Integers only
      \item Tails decay exponentially
    \end{itemize}
    
    \vspace{0.3cm}
    
    \textbf{The Balancing Act:}
    \begin{itemize}
      \item Too small $\sigma$: Insecure
      \item Too large $\sigma$: Decryption fails
      \item Sweet spot: $\sigma \approx \sqrt{n}$
    \end{itemize}
  
    \end{column}
    \begin{column}{0.5\textwidth}

    \begin{tikzpicture}[scale=0.8]
      % Axes
      \draw[thick, ->] (-3,0) -- (3,0) node[right] {$e$};
      \draw[thick, ->] (0,0) -- (0,3.5) node[above] {Probability};
      
      % Discrete points
      \foreach \x in {-5,-4,-3,-2,-1,0,1,2,3,4,5} {
        \pgfmathsetmacro{\prob}{2.5*exp(-\x*\x/4)}
        \draw[blue, ultra thick] (\x/2,0) -- (\x/2,\prob);
        \node[circle, fill=blue, inner sep=1.5pt] at (\x/2,\prob) {};
      }
      
      % Continuous envelope
      \draw[red, dashed, thick, domain=-2.5:2.5, samples=50] 
        plot (\x, {2.5*exp(-4*\x*\x/4)});
      
      % Labels
      \node[below] at (0,0) {$0$};
      \node[below] at (1,0) {$2$};
      \node[below] at (-1,0) {$-2$};
      
      % Annotations
      \draw[<->, green, thick] (-1,1.5) -- (1,1.5) node[midway, above] {$2\sigma$};
    \end{tikzpicture}
    
    \vspace{0.5cm}
    
    \begin{blueBox}
      \textbf{Security comes from:}\\
      Many possible error patterns\\
      $\rightarrow$ Many possible secrets
    \end{blueBox}
  
    \end{column}
  \end{columns}
  
  \teaching{Understand the critical noise component}
\end{frame}

% Slide 3.1.6: Regev's Breakthrough - Worst-Case Hardness
\begin{frame}
  \slideheader{Regev's Breakthrough - Worst-Case Hardness}
  
  \begin{center}
    \textbf{\large The Theoretical Foundation (2005)}
  \end{center}
  
  \vspace{0.3cm}
  
  \begin{redBox}
    \textbf{Regev's Theorem:}\\
    If you can solve LWE on average, you can solve worst-case lattice problems
  \end{redBox}
  
  \vspace{0.3cm}
  
  \begin{tikzpicture}[scale=1]
    % Flow diagram
    \node[draw, thick, minimum width=3cm, minimum height=1cm] (worst) at (0,3) {
      \textbf{ANY lattice instance}
    };
    
    \node[draw, thick, minimum width=3cm, minimum height=1cm, fill=yellow!30] (transform) at (0,1.5) {
      Quantum reduction
    };
    
    \node[draw, thick, minimum width=3cm, minimum height=1cm] (lwe) at (0,0) {
      \textbf{LWE instance}
    };
    
    \node[draw, thick, minimum width=3cm, minimum height=1cm, fill=red!30] (solve) at (4,0) {
      \textbf{LWE solver}
    };
    
    \node[draw, thick, minimum width=3cm, minimum height=1cm] (solution) at (4,3) {
      \textbf{Lattice solution}
    };
    
    % Arrows
    \draw[->, ultra thick] (worst) -- (transform);
    \draw[->, ultra thick] (transform) -- (lwe);
    \draw[->, ultra thick] (lwe) -- (solve);
    \draw[->, ultra thick] (solve) -- (solution);
    \draw[->, ultra thick, dashed] (solution) -- (worst) node[midway, left] {Verify};
  \end{tikzpicture}
  
  \vspace{0.3cm}
  
  \begin{greenBox}
    \textbf{Why This Matters:}
    \begin{itemize}
      \item No ``weak'' LWE instances
      \item Unlike RSA: some moduli might be weak
      \item Average-case = Worst-case security!
    \end{itemize}
  \end{greenBox}
  
  \teaching{Establish strong theoretical foundation}
\end{frame}

% Slide 3.1.7: Building Encryption from LWE
\begin{frame}
  \slideheader{Building Encryption from LWE}
  
  \begin{center}
    \textbf{\large From Hard Problem to Encryption Scheme}
  \end{center}
  
  \vspace{0.3cm}
  
  \begin{blueBox}
    \textbf{Key Generation:}
    \begin{itemize}
      \item Secret key: $\mathbf{s} \in \mathbb{Z}_q^n$
      \item Public key: $(A, \mathbf{b} = A\mathbf{s} + \mathbf{e})$
    \end{itemize}
  \end{blueBox}
  
  \vspace{0.2cm}
  
  \begin{redBox}
    \textbf{Encrypt bit $\mu \in \{0,1\}$:}
    \begin{enumerate}
      \item Sample random subset $S \subseteq \{1,\ldots,m\}$
      \item Compute:
        \begin{align}
          \mathbf{a}' &= \sum_{i \in S} \mathbf{a}_i\\
          b' &= \sum_{i \in S} b_i + \mu \cdot \lfloor q/2 \rfloor
        \end{align}
      \item Ciphertext: $(\mathbf{a}', b')$
    \end{enumerate}
  \end{redBox}
  
  \vspace{0.2cm}
  
  \begin{greenBox}
    \textbf{Decrypt:}
    \begin{itemize}
      \item Compute: $v = b' - \langle \mathbf{a}', \mathbf{s} \rangle$
      \item If $v$ closer to $0$: output $0$
      \item If $v$ closer to $q/2$: output $1$
    \end{itemize}
  \end{greenBox}
  
  \teaching{Show concrete encryption construction}
\end{frame}

% Slide 3.1.8: Why Decryption Works
\begin{frame}
  \slideheader{Why Decryption Works}
  
  \begin{center}
    \textbf{\large Mathematical Analysis of Correctness}
  \end{center}
  
  \vspace{0.3cm}
  
  Let's trace through decryption:
  
  \begin{align}
    v &= b' - \langle \mathbf{a}', \mathbf{s} \rangle\\
    &= \sum_{i \in S} b_i + \mu \cdot \lfloor q/2 \rfloor - \left\langle \sum_{i \in S} \mathbf{a}_i, \mathbf{s} \right\rangle\\
    &= \sum_{i \in S} (\langle \mathbf{a}_i, \mathbf{s} \rangle + e_i) + \mu \cdot \lfloor q/2 \rfloor - \sum_{i \in S} \langle \mathbf{a}_i, \mathbf{s} \rangle\\
    &= \sum_{i \in S} e_i + \mu \cdot \lfloor q/2 \rfloor
  \end{align}
  
  \vspace{0.3cm}
  
  \begin{columns}[T]
    \begin{column}{0.5\textwidth}

    \begin{blueBox}
      \textbf{Key Insight:}
      \begin{itemize}
        \item Error term: $\sum_{i \in S} e_i$ is small
        \item Message term: $\mu \cdot \lfloor q/2 \rfloor$ is large
        \item They don't interfere!
      \end{itemize}
    \end{blueBox}
    
    \vspace{0.3cm}
    
    \textbf{Decryption succeeds if:}
    $$\left|\sum_{i \in S} e_i\right| < \frac{q}{4}$$
    
    High probability for proper parameters!
  
    \end{column}
    \begin{column}{0.5\textwidth}

    \begin{tikzpicture}[scale=0.8]
      % Number line
      \draw[thick, ->] (-0.5,0) -- (5,0) node[right] {$v$};
      
      % Markers
      \draw[thick] (0,-0.2) -- (0,0.2) node[above] {$0$};
      \draw[thick] (2.5,-0.2) -- (2.5,0.2) node[above] {$q/2$};
      \draw[thick] (4.5,-0.2) -- (4.5,0.2) node[above] {$q$};
      
      % Noise clouds
      \draw[blue, thick, fill=blue!20] (-0.5,0.5) -- (0.5,0.5) -- (0.5,1) -- (-0.5,1) -- cycle;
      \node at (0,0.75) {$\mu=0$};
      
      \draw[red, thick, fill=red!20] (2,0.5) -- (3,0.5) -- (3,1) -- (2,1) -- cycle;
      \node at (2.5,0.75) {$\mu=1$};
      
      % Decision boundary
      \draw[dashed, green, thick] (1.25,-0.5) -- (1.25,1.2);
      \node[green, below] at (1.25,-0.6) {Decision};
    \end{tikzpicture}
  
    \end{column}
  \end{columns}
  
  \teaching{Demonstrate correctness mathematically}
\end{frame}

% Slide 3.1.9: From Basic LWE to Ring-LWE
\begin{frame}
  \slideheader{From Basic LWE to Ring-LWE}
  
  \begin{center}
    \textbf{\large The Efficiency Revolution}
  \end{center}
  
  \vspace{0.3cm}
  
  \begin{columns}[T]
    \begin{column}{0.5\textwidth}

    \begin{redBox}
      \textbf{Problem with Basic LWE:}
      \begin{itemize}
        \item Public key: $n \times m$ matrix
        \item Size: $O(n^2)$ elements
        \item For $n=1024$: Several MB!
        \item Too large for practice
      \end{itemize}
    \end{redBox}
    
    \vspace{0.3cm}
    
    \begin{greenBox}
      \textbf{Solution: Add Structure}
      \begin{itemize}
        \item Work in polynomial ring
        \item $R_q = \mathbb{Z}_q[X]/(X^n + 1)$
        \item Structured matrices
        \item 100× size reduction!
      \end{itemize}
    \end{greenBox}
  
    \end{column}
    \begin{column}{0.5\textwidth}

    \begin{tikzpicture}[scale=0.6]
      % Standard LWE matrix
      \node at (0,4) {\textbf{Standard LWE}};
      \draw[thick] (0,3) rectangle (3,0);
      \foreach \x in {0.5,1,1.5,2,2.5} {
        \foreach \y in {0.5,1,1.5,2,2.5} {
          \node at (\x,\y) {\tiny $a_{ij}$};
        }
      }
      \node[below] at (1.5,-0.5) {$n^2$ elements};
      
      % Ring-LWE matrix
      \node at (6,4) {\textbf{Ring-LWE}};
      \draw[thick] (5,3) rectangle (8,0);
      % Circulant structure
      \node at (5.5,2.5) {$a_0$};
      \node at (6.5,2.5) {$a_1$};
      \node at (7.5,2.5) {$a_2$};
      \node at (5.5,1.5) {$a_2$};
      \node at (6.5,1.5) {$a_0$};
      \node at (7.5,1.5) {$a_1$};
      \node at (5.5,0.5) {$a_1$};
      \node at (6.5,0.5) {$a_2$};
      \node at (7.5,0.5) {$a_0$};
      
      \node[below] at (6.5,-0.5) {$n$ elements};
      
      % Arrow showing reduction
      \draw[->, ultra thick, red] (3.5,1.5) -- (4.5,1.5) node[midway, above] {100×};
    \end{tikzpicture}
  
    \end{column}
  \end{columns}
  
  \vspace{0.5cm}
  
  \begin{blueBox}
    \textbf{The Polynomial View:}\\
    Instead of vectors and matrices, work with polynomials:\\
    $a(X) = a_0 + a_1X + a_2X^2 + \cdots + a_{n-1}X^{n-1}$
  \end{blueBox}
  
  \teaching{Motivate practical improvements}
\end{frame}

% Slide 3.1.10: Module-LWE - The Sweet Spot
\begin{frame}
  \slideheader{Module-LWE - The Sweet Spot}
  
  \begin{center}
    \textbf{\large The Spectrum of LWE Variants}
  \end{center}
  
  \vspace{0.5cm}
  
  \begin{tikzpicture}[scale=1.2,
      gradient/.style={
        left color=blue!60,
        right color=red!60,
        middle color=yellow!60
      }
    ]
    % Spectrum bar
    \draw[thick, gradient] (0,0) rectangle (10,1);
    
    % Labels
    \node[above] at (0,1) {\textbf{Plain LWE}};
    \node[above] at (5,1) {\textbf{Module-LWE}};
    \node[above] at (10,1) {\textbf{Ring-LWE}};
    
    % Properties
    \node[below, align=center] at (0,0) {
      \small
      Maximum security\\
      Huge keys\\
      No structure
    };
    
    \node[below, align=center] at (5,0) {
      \small
      Balanced\\
      Moderate keys\\
      Some structure
    };
    
    \node[below, align=center] at (10,0) {
      \small
      Most efficient\\
      Compact keys\\
      Full structure
    };
    
    % NIST choice indicator
    \draw[ultra thick, green, ->] (5,2) -- (5,1.2);
    \node[above, green] at (5,2) {\textbf{NIST Choice}};
  \end{tikzpicture}
  
  \vspace{0.5cm}
  
  \begin{columns}
    \begin{column}{0.33\textwidth}
      \begin{blueBox}
        \textbf{Plain LWE:}\\
        Work with $\mathbb{Z}_q^n$\\
        Full flexibility\\
        Conservative choice
      \end{blueBox}
    \end{column}
    
    \begin{column}{0.33\textwidth}
      \begin{greenBox}
        \textbf{Module-LWE:}\\
        Work with $R_q^k$\\
        $k$ small (2-4)\\
        Best of both worlds
      \end{greenBox}
    \end{column}
    
    \begin{column}{0.33\textwidth}
      \begin{redBox}
        \textbf{Ring-LWE:}\\
        Work with $R_q$\\
        Maximum efficiency\\
        More assumptions
      \end{redBox}
    \end{column}
  \end{columns}
  
  \teaching{Explain the practical choice for standards}
\end{frame}

% ==================================================
% Section 3.2: NIST Standards & Your Future
% ==================================================

\section{NIST Standards \& Your Future}

% Slide 3.2.1: The NIST Competition - Global Effort
\begin{frame}
  \slideheader{The NIST Competition - Global Effort}
  
  \begin{center}
    \textbf{\large The Largest Cryptographic Competition in History}
  \end{center}
  
  \vspace{0.5cm}
  
  \begin{tikzpicture}[scale=1]
    % Timeline
    \draw[ultra thick, ->] (0,0) -- (10,0) node[right] {Time};
    
    % Major milestones
    \draw[thick] (0,-0.2) -- (0,0.2) node[above] {2016};
    \draw[thick] (2,-0.2) -- (2,0.2) node[above] {2017};
    \draw[thick] (4,-0.2) -- (4,0.2) node[above] {2019};
    \draw[thick] (6,-0.2) -- (6,0.2) node[above] {2020};
    \draw[thick] (8,-0.2) -- (8,0.2) node[above] {2022};
    \draw[thick] (10,-0.2) -- (10,0.2) node[above] {2024};
    
    % Events
    \node[below, align=center] at (0,-0.3) {Call for\\Proposals};
    \node[below, align=center] at (2,-0.3) {82\\Submissions};
    \node[below, align=center] at (4,-0.3) {Round 2\\26 candidates};
    \node[below, align=center] at (6,-0.3) {Round 3\\7 finalists};
    \node[below, align=center] at (8,-0.3) {Selection\\4 algorithms};
    \node[below, align=center] at (10,-0.3) {Final\\Standards};
    
    % Stats boxes
    \node[draw, thick, fill=blue!20] at (2,2) {
      \begin{tabular}{c}
        \textbf{82} submissions\\
        \textbf{25} countries\\
        \textbf{1000+} researchers
      \end{tabular}
    };
    
    \node[draw, thick, fill=green!20] at (8,2) {
      \begin{tabular}{c}
        \textbf{5} years\\
        \textbf{3} rounds\\
        \textbf{4} winners
      \end{tabular}
    };
  \end{tikzpicture}
  
  \vspace{0.5cm}
  
  \begin{redBox}
    \textbf{Why This Matters:}
    \begin{itemize}
      \item Global consensus on best algorithms
      \item Extensive cryptanalysis by top experts
      \item Industry and government backing
      \item Your future encryption standard!
    \end{itemize}
  \end{redBox}
  
  \teaching{Show scope and credibility of the process}
\end{frame}

% Slide 3.2.2: The Winners - Lattice Domination
\begin{frame}
  \slideheader{The Winners - Lattice Domination}
  
  \begin{center}
    \textbf{\large NIST Post-Quantum Cryptography Standards}
  \end{center}
  
  \vspace{0.5cm}
  
  \begin{table}
    \centering
    \begin{tabular}{|l|l|l|l|}
      \hline
      \textbf{Algorithm} & \textbf{Type} & \textbf{Based On} & \textbf{Purpose} \\
      \hline
      \cellcolor{green!30}ML-KEM (Kyber) & KEM & Module-LWE & Encryption \\
      \cellcolor{green!30}ML-DSA (Dilithium) & Signature & Module-LWE & Digital Signatures \\
      \cellcolor{green!30}Falcon & Signature & NTRU Lattices & Signatures (alt.) \\
      \cellcolor{yellow!30}SPHINCS+ & Signature & Hash-based & Backup option \\
      \hline
    \end{tabular}
  \end{table}
  
  \vspace{0.5cm}
  
  \begin{columns}
    \begin{column}{0.6\textwidth}
      \begin{tikzpicture}[scale=0.8]
        % Pie chart
        \pie[radius=2]{75/Lattice-based (Green), 25/Hash-based (Yellow)}
      \end{tikzpicture}
    \end{column}
    
    \begin{column}{0.4\textwidth}
      \begin{redBox}
        \textbf{Key Takeaway:}\\
        3 of 4 winners are lattice-based!\\
        ~\\
        Lattices dominate PQC
      \end{redBox}
    \end{column}
  \end{columns}
  
  \vspace{0.3cm}
  
  \begin{greenBox}
    \textbf{Why Lattices Won:}
    Best balance of security, efficiency, and flexibility
  \end{greenBox}
  
  \teaching{Present the new standards landscape}
\end{frame}

% Slide 3.2.3: ML-KEM (Kyber) Overview
\begin{frame}
  \slideheader{ML-KEM (Kyber) Overview}
  
  \begin{center}
    \textbf{\large The New Encryption Standard}
  \end{center}
  
  \vspace{0.3cm}
  
  \begin{columns}[T]
    \begin{column}{0.5\textwidth}

    \textbf{Purpose:} Replace ECDH/RSA for key exchange
    
    \vspace{0.3cm}
    
    \textbf{Key Features:}
    \begin{itemize}
      \item Based on Module-LWE
      \item IND-CCA2 secure
      \item Three security levels
      \item Often faster than ECC!
    \end{itemize}
    
    \vspace{0.3cm}
    
    \textbf{Performance:}
    \begin{itemize}
      \item KeyGen: 12-38 μs
      \item Encaps: 16-48 μs
      \item Decaps: 17-55 μs
    \end{itemize}
  
    \end{column}
    \begin{column}{0.5\textwidth}

    \begin{tabular}{|l|r|r|r|}
      \hline
      \textbf{Level} & \textbf{PK} & \textbf{CT} & \textbf{Security} \\
      \hline
      ML-KEM-512 & 800 B & 768 B & 128-bit \\
      ML-KEM-768 & 1184 B & 1088 B & 192-bit \\
      ML-KEM-1024 & 1568 B & 1568 B & 256-bit \\
      \hline
    \end{tabular}
    
    \vspace{0.5cm}
    
    \begin{tikzpicture}[scale=0.7]
      % Size comparison
      \node[draw, thick, minimum width=0.5cm, minimum height=0.5cm, fill=blue!30] at (0,0) {};
      \node[below] at (0,-0.5) {ECC};
      \node[above] at (0,0.5) {32 B};
      
      \node[draw, thick, minimum width=3cm, minimum height=0.5cm, fill=red!30] at (3,0) {};
      \node[below] at (3,-0.5) {ML-KEM};
      \node[above] at (3,0.5) {1184 B};
      
      \draw[<->, thick] (0.5,1.5) -- (5.5,1.5) node[midway, above] {37× larger};
    \end{tikzpicture}
  
    \end{column}
  \end{columns}
  
  \vspace{0.3cm}
  
  \begin{blueBox}
    \textbf{Already deployed:} Chrome, Cloudflare, AWS, Signal...
  \end{blueBox}
  
  \teaching{Introduce primary encryption standard}
\end{frame}

% Slide 3.2.4: ML-DSA (Dilithium) Overview
\begin{frame}
  \slideheader{ML-DSA (Dilithium) Overview}
  
  \begin{center}
    \textbf{\large The New Signature Standard}
  \end{center}
  
  \vspace{0.3cm}
  
  \begin{columns}[T]
    \begin{column}{0.5\textwidth}

    \textbf{Purpose:} Replace ECDSA/RSA signatures
    
    \vspace{0.3cm}
    
    \textbf{Key Features:}
    \begin{itemize}
      \item ``Fiat-Shamir with Aborts''
      \item Rejection sampling for security
      \item Deterministic signatures
      \item Very fast verification
    \end{itemize}
    
    \vspace{0.3cm}
    
    \textbf{How it Works:}
    \begin{enumerate}
      \item Create candidate signature
      \item Check if it leaks secret
      \item If yes: reject and retry
      \item If no: output signature
    \end{enumerate}
  
    \end{column}
    \begin{column}{0.5\textwidth}

    \begin{tabular}{|l|r|r|r|}
      \hline
      \textbf{Level} & \textbf{PK} & \textbf{Sig} & \textbf{Security} \\
      \hline
      ML-DSA-44 & 1312 B & 2420 B & 128-bit \\
      ML-DSA-65 & 1952 B & 3309 B & 192-bit \\
      ML-DSA-87 & 2592 B & 4627 B & 256-bit \\
      \hline
    \end{tabular}
    
    \vspace{0.5cm}
    
    \textbf{Performance:}
    \begin{itemize}
      \item Sign: 100-300 μs
      \item Verify: 50-150 μs
      \item Faster than RSA verify!
    \end{itemize}
    
    \vspace{0.3cm}
    
    \begin{redBox}
      \textbf{Challenge:}\\
      Signatures 50-100× larger than ECDSA
    \end{redBox}
  
    \end{column}
  \end{columns}
  
  \teaching{Introduce primary signature standard}
\end{frame}

% Slide 3.2.5: Real-World Deployment - Cloudflare
\begin{frame}
  \slideheader{Real-World Deployment - Cloudflare}
  
  \begin{center}
    \textbf{\large Case Study: 20+ Million Websites Protected}
  \end{center}
  
  \vspace{0.5cm}
  
  \begin{columns}[T]
    \begin{column}{0.5\textwidth}

    \textbf{Cloudflare's PQC Timeline:}
    \begin{itemize}
      \item 2019: Initial experiments
      \item 2022: Beta testing with Chrome
      \item 2023: Hybrid mode default
      \item 2024: Full production
    \end{itemize}
    
    \vspace{0.3cm}
    
    \textbf{Deployment Strategy:}
    \begin{itemize}
      \item Hybrid: X25519 + Kyber768
      \item Gradual rollout
      \item Extensive monitoring
      \item Fallback mechanisms
    \end{itemize}
  
    \end{column}
    \begin{column}{0.5\textwidth}

    \begin{tikzpicture}[scale=0.8]
      % Performance impact chart
      \draw[thick, ->] (0,0) -- (4,0) node[right] {Latency};
      \draw[thick, ->] (0,0) -- (0,3) node[above] {\%};
      
      % Bars
      \draw[fill=blue!60] (0.5,0) rectangle (1,2);
      \draw[fill=red!60] (1.5,0) rectangle (2,2.1);
      \draw[fill=green!60] (2.5,0) rectangle (3,0.5);
      
      \node[below] at (0.75,-0.3) {Classic};
      \node[below] at (1.75,-0.3) {Hybrid};
      \node[below] at (2.75,-0.3) {Impact};
      
      \node[above] at (0.75,2) {100\%};
      \node[above] at (1.75,2.1) {105\%};
      \node[above] at (2.75,0.5) {+0.1ms};
    \end{tikzpicture}
    
    \vspace{0.5cm}
    
    \begin{greenBox}
      \textbf{Results:}\\
      $<$0.1ms latency increase\\
      Zero user complaints\\
      Quantum-safe today!
    \end{greenBox}
  
    \end{column}
  \end{columns}
  
  \vspace{0.3cm}
  
  \begin{blueBox}
    \textbf{Lesson:} PQC works at massive scale with minimal impact
  \end{blueBox}
  
  \teaching{Prove practical deployment feasibility}
\end{frame}

% Slide 3.2.6: The Size Challenge
\begin{frame}
  \slideheader{The Size Challenge}
  
  \begin{center}
    \textbf{\large The Elephant in the Room: Key and Signature Sizes}
  \end{center}
  
  \vspace{0.5cm}
  
  \begin{table}
    \centering
    \begin{tabular}{|l|r|r|r|}
      \hline
      \textbf{Algorithm} & \textbf{Public Key} & \textbf{Signature/CT} & \textbf{vs Classic} \\
      \hline
      ECDSA P-256 & 64 B & 64 B & 1× \\
      RSA-2048 & 256 B & 256 B & 4× \\
      \hline
      ML-KEM-768 & 1,184 B & 1,088 B & 18× \\
      ML-DSA-65 & 1,952 B & 3,309 B & 52× \\
      Falcon-512 & 897 B & 666 B & 10× \\
      \hline
    \end{tabular}
  \end{table}
  
  \vspace{0.5cm}
  
  \begin{columns}[T]
    \begin{column}{0.5\textwidth}

    \textbf{Impact Areas:}
    \begin{itemize}
      \item TLS handshakes
      \item Certificate chains
      \item IoT devices
      \item Blockchain systems
      \item Storage requirements
    \end{itemize}
  
    \end{column}
    \begin{column}{0.5\textwidth}

    \textbf{Mitigation Strategies:}
    \begin{itemize}
      \item Caching public keys
      \item Compression techniques
      \item Protocol optimizations
      \item Hybrid modes (transitional)
      \item Hardware acceleration
    \end{itemize}
  
    \end{column}
  \end{columns}
  
  \vspace{0.3cm}
  
  \begin{redBox}
    \textbf{Reality Check:} Size increase is the price of quantum security
  \end{redBox}
  
  \teaching{Address the main practical concern honestly}
\end{frame}

% Slide 3.2.7: Migration Timeline
\begin{frame}
  \slideheader{Migration Timeline}
  
  \begin{center}
    \textbf{\large The Post-Quantum Transition Roadmap}
  \end{center}
  
  \vspace{0.5cm}
  
  \begin{tikzpicture}[scale=0.9]
    % Timeline base
    \draw[ultra thick, ->] (0,0) -- (10,0) node[right] {Year};
    
    % Year markers
    \foreach \year/\x in {2024/0, 2025/2, 2027/4, 2030/6, 2035/8, 2040/10} {
      \draw[thick] (\x,-0.2) -- (\x,0.2) node[above] {\year};
    }
    
    % Phases
    \draw[fill=green!30, draw=green, thick] (0,1) rectangle (2,2);
    \node at (1,1.5) {Standards};
    
    \draw[fill=blue!30, draw=blue, thick] (2,1) rectangle (4,2);
    \node at (3,1.5) {Early Adopters};
    
    \draw[fill=yellow!30, draw=yellow!80!black, thick] (4,1) rectangle (6,2);
    \node at (5,1.5) {Mainstream};
    
    \draw[fill=red!30, draw=red, thick] (6,1) rectangle (8,2);
    \node at (7,1.5) {Legacy Sunset};
    
    % Danger zone
    \draw[ultra thick, red, dashed] (8,0) -- (8,3);
    \node[red, right, align=left] at (8.2,2.5) {Quantum\\Threat\\Real};
    
    % Harvest now, decrypt later
    \draw[<->, thick, purple] (0,-1.5) -- (8,-1.5);
    \node[purple, below] at (4,-1.7) {``Harvest Now, Decrypt Later'' Risk};
  \end{tikzpicture}
  
  \vspace{0.5cm}
  
  \begin{columns}
    \begin{column}{0.5\textwidth}
      \begin{blueBox}
        \textbf{NOW (2024-2025):}
        \begin{itemize}
          \item Inventory systems
          \item Test implementations
          \item Plan migration
        \end{itemize}
      \end{blueBox}
    \end{column}
    
    \begin{column}{0.5\textwidth}
      \begin{redBox}
        \textbf{URGENT:}
        \begin{itemize}
          \item Long-term secrets
          \item Critical infrastructure
          \item Financial systems
        \end{itemize}
      \end{redBox}
    \end{column}
  \end{columns}
  
  \teaching{Create urgency for action}
\end{frame}

% Slide 3.2.8: Hybrid Modes - Belt and Suspenders
\begin{frame}
  \slideheader{Hybrid Modes - Belt and Suspenders}
  
  \begin{center}
    \textbf{\large Safe Transition Strategy: Use Both!}
  \end{center}
  
  \vspace{0.5cm}
  
  \begin{tikzpicture}[scale=1]
    % Classical crypto
    \node[draw, thick, minimum width=2.5cm, minimum height=1.5cm, fill=blue!20] (classical) at (-3,0) {
      \textbf{Classical}\\
      (ECDH)
    };
    
    % PQC
    \node[draw, thick, minimum width=2.5cm, minimum height=1.5cm, fill=green!20] (pqc) at (3,0) {
      \textbf{Post-Quantum}\\
      (ML-KEM)
    };
    
    % Combination
    \node[draw, ultra thick, minimum width=2cm, minimum height=1cm, fill=yellow!30] (combine) at (0,-2) {
      \textbf{Combine}\\
      Keys
    };
    
    % Arrows
    \draw[->, thick] (classical) -- (combine);
    \draw[->, thick] (pqc) -- (combine);
    
    % Result
    \node[draw, ultra thick, minimum width=3cm, minimum height=1cm, fill=red!30] (result) at (0,-4) {
      \textbf{Hybrid Security}
    };
    
    \draw[->, ultra thick] (combine) -- (result);
    
    % Annotations
    \node[right] at (4,0) {
      \begin{tabular}{l}
        ✓ Proven today\\
        ✗ Quantum vulnerable
      \end{tabular}
    };
    
    \node[left] at (-4,0) {
      \begin{tabular}{l}
        ✓ Quantum safe\\
        ? New, less tested
      \end{tabular}
    };
    
    \node[below] at (0,-5) {
      \textbf{Must break BOTH to compromise!}
    };
  \end{tikzpicture}
  
  \vspace{0.3cm}
  
  \begin{greenBox}
    \textbf{Recommended during transition:}
    Best of both worlds - no security regression
  \end{greenBox}
  
  \teaching{Explain safe migration approach}
\end{frame}

% Slide 3.2.9: Your Unique Advantages as Physicists
\begin{frame}
  \slideheader{Your Unique Advantages as Physicists}
  
  \begin{center}
    \textbf{\large Why YOU Are Perfectly Positioned}
  \end{center}
  
  \vspace{0.5cm}
  
  \begin{columns}[T]
    \begin{column}{0.5\textwidth}

    \begin{greenBox}
      \textbf{Technical Advantages:}
      \begin{itemize}
        \item[\checkmark] Real quantum intuition
        \item[\checkmark] Lattice experience (crystals)
        \item[\checkmark] Noise/error understanding
        \item[\checkmark] Hardware knowledge
        \item[\checkmark] Mathematical maturity
      \end{itemize}
    \end{greenBox}
    
    \vspace{0.3cm}
    
    \begin{greenBox}
      \textbf{Perspective Advantages:}
      \begin{itemize}
        \item[\checkmark] Bridge theory \& experiment
        \item[\checkmark] Systems thinking
        \item[\checkmark] Physical constraints focus
        \item[\checkmark] Different viewpoint from CS
      \end{itemize}
    \end{greenBox}
  
    \end{column}
    \begin{column}{0.5\textwidth}

    \begin{tikzpicture}[scale=0.8]
      % Venn diagram
      \draw[thick, fill=blue!20] (0,0) circle (2cm);
      \draw[thick, fill=red!20] (2.5,0) circle (2cm);
      \draw[thick, fill=green!20] (1.25,-2) circle (2cm);
      
      \node at (-0.5,1) {\small \textbf{Quantum}};
      \node at (3,1) {\small \textbf{Crypto}};
      \node at (1.25,-3) {\small \textbf{Physics}};
      
      \node[font=\large\bfseries] at (1.25,0) {YOU!};
      
      % Arrows pointing to opportunities
      \draw[->, ultra thick] (3.5,2) -- (2,0.5);
      \node[right] at (3.5,2) {Research};
      
      \draw[->, ultra thick] (3.5,-2) -- (2,-0.5);
      \node[right] at (3.5,-2) {Industry};
      
      \draw[->, ultra thick] (-2,2) -- (0.5,0.5);
      \node[left] at (-2,2) {Academia};
    \end{tikzpicture}
  
    \end{column}
  \end{columns}
  
  \vspace{0.3cm}
  
  \begin{redBox}
    \textbf{The Field Needs You:}
    Your physics background provides insights that pure CS/Math people miss
  \end{redBox}
  
  \teaching{Reinforce their unique qualifications}
\end{frame}

% Slide 3.2.10: Research Opportunities
\begin{frame}
  \slideheader{Research Opportunities}
  
  \begin{center}
    \textbf{\large Open Problems Where Physics Meets PQC}
  \end{center}
  
  \vspace{0.5cm}
  
  \begin{columns}
    \begin{column}{0.5\textwidth}
      \textbf{Theoretical Directions:}
      \begin{itemize}
        \item Quantum algorithms for lattices
        \item New hardness assumptions
        \item Connections to condensed matter
        \item Quantum reduction improvements
        \item Novel cryptographic constructions
      \end{itemize}
      
      \vspace{0.3cm}
      
      \textbf{Physical Perspectives:}
      \begin{itemize}
        \item Thermodynamic crypto bounds
        \item Photonic implementations
        \item Quantum error analogies
        \item Phase transition insights
      \end{itemize}
    \end{column}
    
    \begin{column}{0.5\textwidth}
      \textbf{Practical Challenges:}
      \begin{itemize}
        \item Side-channel resistance
        \item Hardware acceleration
        \item Quantum RNG design
        \item Hybrid protocol optimization
        \item IoT implementations
      \end{itemize}
      
      \vspace{0.3cm}
      
      \textbf{Interdisciplinary:}
      \begin{itemize}
        \item Quantum-classical interfaces
        \item Error correction crossover
        \item Physical unclonable functions
        \item Quantum key distribution + PQC
      \end{itemize}
    \end{column}
  \end{columns}
  
  \vspace{0.5cm}
  
  \begin{blueBox}
    \textbf{Your PhD could:} Pioneer new connections between quantum physics and post-quantum cryptography
  \end{blueBox}
  
  \teaching{Inspire research directions}
\end{frame}

% Slide 3.2.11: Practical Next Steps
\begin{frame}
  \slideheader{Practical Next Steps}
  
  \begin{center}
    \textbf{\large Your Action Plan}
  \end{center}
  
  \vspace{0.5cm}
  
  \begin{tikzpicture}[scale=1]
    % Timeline
    \draw[ultra thick, ->] (0,0) -- (10,0);
    
    % This week
    \node[draw, thick, fill=green!30, minimum width=2cm, minimum height=1.5cm] at (1,2) {
      \begin{tabular}{c}
        \textbf{This Week}\\
        Download NIST\\
        standards\\
        Try liboqs\\
        Run demos
      \end{tabular}
    };
    \draw[->, thick] (1,1) -- (1,0.2);
    
    % This month
    \node[draw, thick, fill=blue!30, minimum width=2cm, minimum height=1.5cm] at (4,2) {
      \begin{tabular}{c}
        \textbf{This Month}\\
        Implement\\
        baby Kyber\\
        Join forums\\
        Read papers
      \end{tabular}
    };
    \draw[->, thick] (4,1) -- (4,0.2);
    
    % This year
    \node[draw, thick, fill=yellow!30, minimum width=2cm, minimum height=1.5cm] at (7,2) {
      \begin{tabular}{c}
        \textbf{This Year}\\
        Research\\
        project\\
        Conference\\
        Contribute
      \end{tabular}
    };
    \draw[->, thick] (7,1) -- (7,0.2);
    
    % Career
    \node[draw, thick, fill=red!30, minimum width=2cm, minimum height=1.5cm] at (10,2) {
      \begin{tabular}{c}
        \textbf{Career}\\
        Lead PQC\\
        transition\\
        Shape the\\
        future!
      \end{tabular}
    };
    \draw[->, thick] (10,1) -- (10,0.2);
  \end{tikzpicture}
  
  \vspace{0.5cm}
  
  \begin{greenBox}
    \textbf{Start Today:}
    \url{https://github.com/open-quantum-safe/liboqs}
  \end{greenBox}
  
  \teaching{Concrete actionable steps}
\end{frame}

% Slide 3.2.12: Resources for Deep Dive
\begin{frame}
  \slideheader{Resources for Deep Dive}
  
  \begin{center}
    \textbf{\large Your Learning Toolkit}
  \end{center}
  
  \vspace{0.5cm}
  
  \begin{columns}[T]
    \begin{column}{0.5\textwidth}

    \textbf{Essential Papers:}
    \begin{itemize}
      \item Regev (2005) - LWE introduction
      \item Peikert (2016) - ``A Decade of Lattice''
      \item NIST standards (2024)
      \item Micciancio-Regev survey
    \end{itemize}
    
    \vspace{0.3cm}
    
    \textbf{Online Courses:}
    \begin{itemize}
      \item Coursera - Post-Quantum Crypto
      \item MIT OCW - Lattice Cryptography
      \item YouTube - Simons Institute
    \end{itemize}
  
    \end{column}
    \begin{column}{0.5\textwidth}

    \textbf{Code Libraries:}
    \begin{itemize}
      \item \texttt{liboqs} - All NIST algorithms
      \item \texttt{PQClean} - Reference implementations
      \item \texttt{CRYSTALS} - Kyber/Dilithium
      \item \texttt{FrodoKEM} - Learning project
    \end{itemize}
    
    \vspace{0.3cm}
    
    \textbf{Communities:}
    \begin{itemize}
      \item NIST PQC Forum
      \item IACR ePrint Archive
      \item PQC Standardization list
      \item Stack Exchange - Crypto
    \end{itemize}
  
    \end{column}
  \end{columns}
  
  \vspace{0.5cm}
  
  \begin{blueBox}
    \textbf{Pro Tip:} Start with Peikert's survey, then implement toy LWE encryption
  \end{blueBox}
  
  \teaching{Enable continued self-learning}
\end{frame}

% Slide 3.2.13: The Big Picture
\begin{frame}
  \slideheader{The Big Picture}
  
  \begin{center}
    \textbf{\Huge Remember These Four Things}
  \end{center}
  
  \vspace{1cm}
  
  \begin{enumerate}
    \item \begin{redBox}
      \textbf{Current crypto WILL fail against quantum computers}\\
      RSA, ECDSA, DH - all broken by Shor's algorithm
    \end{redBox}
    
    \vspace{0.5cm}
    
    \item \begin{greenBox}
      \textbf{Lattice-based solutions are ready NOW}\\
      NIST standards published, real-world deployments working
    \end{greenBox}
    
    \vspace{0.5cm}
    
    \item \begin{blueBox}
      \textbf{Migration must start IMMEDIATELY}\\
      ``Harvest now, decrypt later'' means secrets already at risk
    \end{blueBox}
    
    \vspace{0.5cm}
    
    \item \begin{yellowBox}
      \textbf{YOU can make unique contributions}\\
      Physics background + quantum knowledge = perfect combination
    \end{yellowBox}
  \end{enumerate}
  
  \teaching{Reinforce the four key messages}
\end{frame}

% Slide 3.2.14: Your Mission
\begin{frame}
  \slideheader{Your Mission}
  
  \vfill
  
  \begin{center}
    \begin{tikzpicture}[scale=1.2]
      % Quote box
      \node[draw, ultra thick, minimum width=10cm, minimum height=4cm, fill=blue!10] at (0,0) {
        \begin{minipage}{9cm}
          \centering
          \Large
          \textit{``In 2074, historians will write about the post-quantum transition.}
          
          \vspace{0.5cm}
          
          \textit{Will they write:}
          
          \vspace{0.3cm}
          
          \textit{`Scientists acted in time'}
          
          \textit{or}
          
          \textit{`If only they had started sooner'?''}
          
          \vspace{0.5cm}
          
          \normalsize
          \textbf{That story is being written by YOUR actions.}
        \end{minipage}
      };
    \end{tikzpicture}
  \end{center}
  
  \vfill
  
  \begin{center}
    \Huge \textbf{The future needs you. Act now.}
  \end{center}
  
  \vfill
  
  \teaching{Inspirational call to action}
\end{frame}

% Slide 3.2.15: Questions & Discussion
\begin{frame}
  \slideheader{Questions \& Discussion}
  
  \begin{center}
    \textbf{\Large Thank you for your attention!}
  \end{center}
  
  \vspace{1cm}
  
  \begin{columns}
    \begin{column}{0.5\textwidth}
      \begin{blueBox}
        \textbf{Contact Information:}\\
        Dr. [Your Name]\\
        [Your Email]\\
        [Your Institution]\\
        Office Hours: [Times]
      \end{blueBox}
      
      \vspace{0.5cm}
      
      \begin{greenBox}
        \textbf{Research Opportunities:}\\
        - PhD positions available\\
        - Summer research projects\\
        - Industry collaborations
      \end{greenBox}
    \end{column}
    
    \begin{column}{0.5\textwidth}
      \begin{center}
        \begin{tikzpicture}[scale=0.8]
          % Question marks in creative arrangement
          \node[font=\Huge, rotate=20, color=blue] at (0,2) {?};
          \node[font=\huge, rotate=-15, color=red] at (1.5,1) {?};
          \node[font=\Large, rotate=35, color=green] at (-1,1) {?};
          \node[font=\LARGE, rotate=-25, color=purple] at (0.5,0) {?};
          \node[font=\large, rotate=10, color=orange] at (-0.5,2.5) {?};
        \end{tikzpicture}
      \end{center}
      
      \vspace{0.5cm}
      
      \textbf{Discussion Starters:}
      \begin{itemize}
        \item What excites you most?
        \item What concerns you?
        \item How does this connect to your research?
        \item Where do you see opportunities?
      \end{itemize}
    \end{column}
  \end{columns}
  
  \vspace{0.5cm}
  
  \begin{center}
    \large
    \textbf{Remember:} You're not just learning about the future of cryptography\\
    --- you're going to help build it!
  \end{center}
  
  \teaching{Open dialogue, gauge understanding, inspire action}
\end{frame}

% ==================================================
% BACKUP/OVERFLOW SLIDES
% ==================================================

\section*{Backup Slides}

% Slide B.1: Why Not Just Bigger RSA Keys?
\begin{frame}
  \slideheader{Why Not Just Bigger RSA Keys?}
  
  \begin{center}
    \textbf{\large Common Question: Can't We Just Use RSA-1000000?}
  \end{center}
  
  \vspace{0.5cm}
  
  \begin{columns}[T]
    \begin{column}{0.5\textwidth}

    \textbf{The Fundamental Problem:}
    
    \begin{itemize}
      \item Classical: $\exp(n^{1/3})$ time
      \item Quantum: $n^3$ time
      \item Gap: $\frac{\exp(n^{1/3})}{n^3}$
    \end{itemize}
    
    \vspace{0.3cm}
    
    \textbf{No key size saves you!}
    
    \begin{tabular}{|l|r|r|}
      \hline
      \textbf{RSA Size} & \textbf{Classical} & \textbf{Quantum} \\
      \hline
      2048 bits & $10^{20}$ years & 8 hours \\
      10,000 bits & $10^{100}$ years & 1 day \\
      1M bits & $10^{10000}$ years & 1 year \\
      \hline
    \end{tabular}
  
    \end{column}
    \begin{column}{0.5\textwidth}

    \begin{tikzpicture}[scale=0.7]
      % Log-log plot
      \draw[thick, ->] (0,0) -- (5,0) node[right] {Key size (log)};
      \draw[thick, ->] (0,0) -- (0,4) node[above] {Time (log)};
      
      % Classical exponential
      \draw[red, ultra thick, domain=0.5:4.5] plot (\x, {0.5*exp(0.3*\x)});
      \node[red, right] at (4.5,3.8) {Classical};
      
      % Quantum polynomial
      \draw[blue, ultra thick] plot[smooth] coordinates {
        (0.5,0.2) (1,0.3) (2,0.5) (3,0.7) (4,0.9) (4.5,1)
      };
      \node[blue, right] at (4.5,1) {Quantum};
      
      % Gap grows
      \draw[<->, green, thick] (3,0.7) -- (3,3.2);
      \node[green, right] at (3,2) {Gap grows!};
    \end{tikzpicture}
    
    \vspace{0.3cm}
    
    \textbf{Plus practical issues:}
    \begin{itemize}
      \item GB-sized keys
      \item Minutes per operation
      \item Completely unusable
    \end{itemize}
  
    \end{column}
  \end{columns}
  
  \vspace{0.3cm}
  
  \begin{redBox}
    \textbf{Answer:} No amount of key size increase can overcome exponential vs polynomial gap
  \end{redBox}
  
  \teaching{Address common misconception}
\end{frame}

% Slide B.2: Detailed Shor Circuit
\begin{frame}
  \slideheader{Detailed Shor Circuit}
  
  \begin{center}
    \textbf{\large For Those Who Want the Quantum Details}
  \end{center}
  
  \vspace{0.3cm}
  
  \begin{tikzpicture}[scale=0.8]
    % Quantum circuit
    % Register 1
    \node[left] at (-1,4) {$|0\rangle$};
    \draw[thick] (0,4) -- (11,4);
    
    \node[left] at (-1,3.5) {$|0\rangle$};
    \draw[thick] (0,3.5) -- (11,3.5);
    
    \node at (0,3) {$\vdots$};
    
    \node[left] at (-1,2.5) {$|0\rangle$};
    \draw[thick] (0,2.5) -- (11,2.5);
    
    % Register 2
    \draw[thick, double] (0,1.5) -- (11,1.5);
    \node[left] at (-1,1.5) {$|0\rangle^{\otimes n}$};
    
    % Hadamards
    \foreach \y in {4,3.5,2.5} {
      \draw[thick, fill=blue!20] (0.5,\y-0.2) rectangle (1,\y+0.2);
      \node at (0.75,\y) {$H$};
    }
    
    % Controlled modular exponentiation
    \draw[thick, fill=red!20] (2,2) rectangle (5,4.5);
    \node[align=center] at (3.5,3.25) {Controlled\\$U^{2^j}$};
    
    % Control dots
    \foreach \y in {4,3.5,2.5} {
      \node[circle, fill=black, inner sep=1.5pt] at (3.5,\y) {};
      \draw[thick] (3.5,\y) -- (3.5,2);
    }
    
    % QFT inverse
    \draw[thick, fill=green!20] (6,2.3) rectangle (8,4.5);
    \node at (7,3.4) {QFT$^{-1}$};
    
    % Measurements
    \foreach \y in {4,3.5,2.5} {
      \draw[thick] (9,\y-0.2) -- (9,\y+0.2) -- (9.5,\y+0.2) -- (9.5,\y-0.2) -- (9,\y-0.2);
      \draw[thick] (9.25,\y) -- (9.25,\y-0.3) -- (9.5,\y-0.5);
    }
    
    % Classical processing
    \draw[thick, dashed, ->] (10,3.5) -- (12,3.5) node[right] {Classical: extract period};
  \end{tikzpicture}
  
  \vspace{0.5cm}
  
  \textbf{Gate Counts (for RSA-2048):}
  \begin{itemize}
    \item Total gates: $\sim 10^9$
    \item Depth: $\sim 10^8$
    \item Logical qubits: $\sim 4000$
    \item T gates (costly): $\sim 10^8$
  \end{itemize}
  
  \vspace{0.3cm}
  
  \begin{blueBox}
    \textbf{Key Resource:} Modular exponentiation dominates cost ($>99\%$ of circuit)
  \end{blueBox}
  
  \teaching{Technical deep dive for interested students}
\end{frame}

% Slide B.3: LWE Security Proof Sketch
\begin{frame}
  \slideheader{LWE Security Proof Sketch}
  
  \begin{center}
    \textbf{\large High-Level Idea of Regev's Reduction}
  \end{center}
  
  \vspace{0.5cm}
  
  \textbf{Goal:} Show that solving average-case LWE $\Rightarrow$ solving worst-case lattice problems
  
  \vspace{0.3cm}
  
  \begin{enumerate}
    \item \textbf{Start with worst-case lattice:}
    \begin{itemize}
      \item Given: Arbitrary lattice $\mathcal{L}$
      \item Want: Short vector in $\mathcal{L}$
    \end{itemize}
    
    \item \textbf{Quantum step - create superposition:}
    \begin{align}
      |\psi\rangle = \sum_{\mathbf{v} \in \mathcal{L}} \rho(||\mathbf{v}||) |\mathbf{v}\rangle
    \end{align}
    where $\rho$ is Gaussian weight
    
    \item \textbf{Measure to get LWE samples:}
    \begin{itemize}
      \item Measurement produces LWE distribution
      \item Secret related to dual lattice
    \end{itemize}
    
    \item \textbf{Use LWE solver:}
    \begin{itemize}
      \item Feed samples to assumed LWE solver
      \item Get information about dual lattice
    \end{itemize}
    
    \item \textbf{Recover short vector:}
    \begin{itemize}
      \item Use dual information to find short primal vector
      \item Quantum amplitude amplification helps
    \end{itemize}
  \end{enumerate}
  
  \vspace{0.3cm}
  
  \begin{greenBox}
    \textbf{Key Insight:} LWE samples are ``quantum snapshots'' of the lattice
  \end{greenBox}
  
  \teaching{High-level theoretical understanding}
\end{frame}

% Slide B.4: Implementation Code Example
\begin{frame}[fragile]
  \slideheader{Implementation Code Example}
  
  \begin{center}
    \textbf{\large Simple Python LWE Encryption}
  \end{center}
  
  \vspace{0.3cm}
  
  \begin{lstlisting}[language=Python, basicstyle=\tiny]
import numpy as np

class ToyLWE:
    def __init__(self, n=10, q=97, sigma=3):
        self.n, self.q, self.sigma = n, q, sigma
        
    def keygen(self):
        # Secret key
        self.s = np.random.randint(0, self.q, self.n)
        
        # Public key: (A, b = As + e)
        self.A = np.random.randint(0, self.q, (self.n, self.n))
        e = np.random.normal(0, self.sigma, self.n).astype(int)
        self.b = (self.A @ self.s + e) % self.q
        
        return (self.A, self.b), self.s
    
    def encrypt(self, bit):
        # Random subset
        r = np.random.randint(0, 2, self.n)
        
        # Encryption
        a_prime = (r @ self.A) % self.q
        b_prime = (r @ self.b + bit * (self.q // 2)) % self.q
        
        return (a_prime, b_prime)
    
    def decrypt(self, ct, sk):
        a_prime, b_prime = ct
        v = (b_prime - np.dot(a_prime, sk)) % self.q
        
        # Closer to 0 or q/2?
        if v < self.q // 4 or v > 3 * self.q // 4:
            return 0
        else:
            return 1

# Test it!
lwe = ToyLWE()
pk, sk = lwe.keygen()
ct = lwe.encrypt(1)
print(f"Decrypted: {lwe.decrypt(ct, sk)}")  # Should print 1
  \end{lstlisting}
  
  \vspace{0.3cm}
  
  \begin{blueBox}
    \textbf{Try during break:} Modify parameters and see when decryption fails!
  \end{blueBox}
  
  \begin{greenBox}
    \textbf{QR Code to full repository:} [QR code would go here]
  \end{greenBox}
  
  \teaching{Hands-on learning opportunity}
\end{frame}

% Slide B.5: Current PQC Adoption Status
\begin{frame}
  \slideheader{Current PQC Adoption Status}
  
  \begin{center}
    \textbf{\large Who's Already Using Post-Quantum Crypto?}
  \end{center}
  
  \vspace{0.5cm}
  
  \begin{table}
    \centering
    \small
    \begin{tabular}{|l|l|l|l|}
      \hline
      \textbf{Company} & \textbf{Product} & \textbf{Algorithm} & \textbf{Status} \\
      \hline
      Google & Chrome & Kyber & Default 2024 \\
      Cloudflare & CDN & Kyber & Production \\
      AWS & S3, KMS & Kyber & Available \\
      Microsoft & Azure & Multiple & Preview \\
      Signal & Messaging & Kyber & Deployed \\
      Apple & iMessage & Kyber & Announced \\
      Zoom & Video & Kyber & Testing \\
      \hline
    \end{tabular}
  \end{table}
  
  \vspace{0.5cm}
  
  \begin{columns}[T]
    \begin{column}{0.5\textwidth}

    \textbf{Government Initiatives:}
    \begin{itemize}
      \item US: NSA mandate 2025
      \item EU: Quantum flagship
      \item China: Massive investment
      \item UK: National strategy
    \end{itemize}
  
    \end{column}
    \begin{column}{0.5\textwidth}

    \textbf{Success Metrics:}
    \begin{itemize}
      \item Billions of PQC handshakes daily
      \item No significant issues reported
      \item Performance acceptable
      \item Transparent to users
    \end{itemize}
  
    \end{column}
  \end{columns}
  
  \vspace{0.3cm}
  
  \begin{greenBox}
    \textbf{The transition is happening NOW - be part of it!}
  \end{greenBox}
  
  \teaching{Show real-world momentum}
\end{frame}

% ==================================================
% END OF PRESENTATION
% ==================================================

\end{document}
